% The fourth order at every rank.
% Alexander Smith, 9 September 2026.
% Companion repository: WORKHOUSE. Every marked statement is machine-checked;
% see Section "Evidence" and the coverage appendix.
\documentclass[11pt]{article}

\usepackage[margin=1in]{geometry}
\usepackage{amsmath,amssymb,amsthm}
\usepackage{mathtools}
\usepackage{booktabs}
\usepackage{array}
\usepackage[hidelinks]{hyperref}
\usepackage{microtype}

\newtheorem{theorem}{Theorem}[section]
\newtheorem{proposition}[theorem]{Proposition}
\newtheorem{lemma}[theorem]{Lemma}
\newtheorem{corollary}[theorem]{Corollary}
\theoremstyle{definition}
\newtheorem{definition}[theorem]{Definition}
\newtheorem{hypothesis}[theorem]{Hypothesis}
\theoremstyle{remark}
\newtheorem{remark}[theorem]{Remark}
\newtheorem{finding}[theorem]{Finding}

\DeclareMathOperator{\spec}{spec}
\DeclareMathOperator{\Tr}{Tr}
\DeclareMathOperator{\adj}{Adj}
\DeclareMathOperator{\Sym}{Sym}
\DeclareMathOperator{\rank}{rank}
\newcommand{\CF}{C_F}
\newcommand{\dg}{^{\dagger}}
\newcommand{\Ldn}{L_{\downarrow}}
\newcommand{\Lup}{L_{\uparrow}}
\newcommand{\Ssq}{S_{\square}}
\newcommand{\Q}{\mathbb{Q}}
\newcommand{\Z}{\mathbb{Z}}
\newcommand{\C}{\mathbb{C}}

% The marker below names, verbatim, a registered machine check that carries the
% sentence it follows. It typesets nothing: it is read by the repository's
% coverage guard, which fails if the named check does not exist or does not
% pass. Same convention, and same macro body, as the master edition.
\newcommand{\chk}[1]{}

\title{The fourth order at every rank:\\
closed forms, a Hodge decomposition, and the mechanism of the\\
tier collapse in strong-coupling $SU(N)$ Hamiltonian lattice gauge theory}

\author{Alexander Smith}
\date{9 September 2026}

\begin{document}
\maketitle

\begin{abstract}
\noindent
We study the Kogut--Susskind Hamiltonian for pure $SU(N)$ gauge theory on a
periodic cubic spatial lattice at strong coupling, in the normalised
charge-odd fundamental-plaquette sector and in the all-rank coordinate
$u=\beta_N/(2N)$. The second-order effective operator in that sector is
$t_Nu^2BB\dg$ with $B=\partial_2\dg$ the oriented cellular coboundary and
$t_N=2N(N^2-4)/[(N^2-1)(2N^2-1)(4N^2-9)]$; its carrier is the homology-protected
flat band $\ker\partial_2$. This paper settles the fourth order on that carrier.

Four results. \emph{First}, a structure theorem: the fourth-order kernel lies
in the incidence algebra generated by the two plaquette Laplacians
$\Ldn=\partial_2\dg\partial_2$ and $\Lup=\partial_3\partial_3\dg$ up to exactly
one operator $R$, the cross-plane half of the signed shared-link adjacency,
\[
  H_4=-\tilde\nu(\Lup-2)+u\,\Ssq^2-\tilde\pi\,\Ssq+\tilde\sigma\,I-2C_{\mathrm{shp}}R,
  \qquad \tilde\nu=-\tfrac{5}{48},
\]
exactly on all $189$ records of the historical kernel. \emph{Second}, a
mechanism for the ``tier collapse'' $B_{\mathrm{shp}}=D_{\mathrm{shp}}=0$, the
vanishing of the two top-degree shape coefficients that a vertex-counting
degree bound was once --- wrongly --- thought to forbid. At every Bloch point
with $q>0$ the Feshbach complement of the carrier is exactly $\ker\Lup$, so
neither Hodge generator can leave the retained sector and every excursion costs
one insertion of $R$; the actual fourth-order support $\{I,\Ssq,\Ssq^2,\Lup,R\}$
has carrier symbols in $\operatorname{span}\{q,q^2,e_2\}$, which is the
collapse. The remaining dynamical input --- that one weight $u$ multiplies every
two-hop chain geometry --- is reduced to a finite identity of Haar integrals and
proved by exhaustion on the fourth-order family.

\emph{Third}, closed forms. Running a from-scratch Wilson-loop calculus over the
rational function field $\Q(N)$ rather than rank by rank, every cluster cumulant
of the fourth order but one is \emph{derived} as a single rational function of
$N$: the two-hop weight, both dressings of both pairs, the corner, and both cube
completions. The exception is the pair cluster, which is never computed and does
not need to be, because a parity argument makes it cancel identically. The
assembly then collapses to three cumulants plus the adjacent-face cube
completion $-106/[N(N^2-1)^3]$, and
\[
  \beta_N=-16u+32d_{\mathrm{single}}-16d_{\mathrm{corner}}+\frac{848}{N(N^2-1)^3}
        =\frac{P_{17}(N^2)}{N\,R_{20}(N^2)},
\]
a polynomial identity of degree $34$ in $N$, proof-checked in Lean~4. This
removes the degree hypothesis under which the corresponding rank-sweep
identification was previously only conditional. The derivation argument reaches
$N\ge5$; at $N=3$ and $N=4$ the forms agree with the per-rank engine as a
checked fact.

\emph{Fourth}, the fourth-order off-axis coefficient, the program's last open
contradiction, is adjudicated --- against both previously recorded kernels. Solving
the shape ansatz over the whole Brillouin zone rather than fitting it at four
points gives $A=5/48$, $B=D=0$ and no residual; the disagreement reduces to the
single scalar $C_{\mathrm{shp}}=-5/96-u-(\rho+\tilde\pi)/2$; and a third,
independent implementation of the corner cluster together with the original
pipeline's own recovered word ledger localises the error to sixteen missing
adjacent-face cube insertion orderings, so
$C_{\mathrm{shp}}=C_{\mathrm{historical}}+25/1024
=-13035490122347/550663802582400$. The $SU(3)$ shape row moves accordingly:
$\beta_3$, $W_4$ and $\lambda_R$ by $25/64$, $\lambda_M$ by $25/128$, with
$\alpha$ and $\lambda_X$ unchanged.

We are explicit about what is \emph{not} established. Every statement here is a
fixed-lattice-spacing statement about an effective operator on the projected
plaquette manifold, downstream of four premises stated in
Section~\ref{sec:premises}. No continuum limit, no physical glueball mass and
no infinite-volume mass gap is claimed; two of the obstructions to the obvious
route are themselves theorems. The off-axis result is an adjudication whose own
declared falsifier --- a direct balanced contraction at $N=3$ --- has not been
performed, the universality lemma is exhausted at three of its four
fourth-order multiplicities, every all-rank statement flows through one
implementation, and no external group has reproduced anything in this program.
Section~\ref{sec:scope} lists each of these with the experiment that would
settle it.
\end{abstract}

\tableofcontents

\section{Introduction}\label{sec:intro}

\subsection{The object}

Fix $N\ge3$ and a periodic cubic spatial lattice $\Lambda=(\Z/L\Z)^3$, $L\ge3$.
The Kogut--Susskind Hamiltonian for pure $SU(N)$ gauge theory is
\begin{equation}\label{eq:ks}
  H \;=\; H_E \;-\; u\,V ,\qquad
  H_E=\sum_{e}\,\tfrac12\sum_a (E_e^a)^2,\qquad
  V=\sum_{p}\bigl(\Tr W_p+\Tr W_p\dg\bigr),
\end{equation}
in the all-rank strong-coupling coordinate $u=\beta_N/(2N)$. (The $SU(3)$
specialisation $u=\beta_3/6$ is the coordinate most of the older literature
uses; the archived $Y=4u$ line is a label erratum and no $4^r$ rescaling of
coefficients is permitted.) The electric term is diagonal in the flux basis and
the magnetic term moves one plaquette at a time.

Degenerate perturbation theory about the electric vacuum, restricted to the
\emph{charge-odd, trivial-flux, fundamental one-plaquette shell} --- the
$3L^3$-dimensional space spanned by single plaquettes carrying fundamental flux,
antisymmetrised in charge conjugation --- produces an effective Hamiltonian
\begin{equation}\label{eq:eff}
  H_{\mathrm{eff}}(u)\;=\;\text{const}\;+\;u^2H_2\;+\;u^3H_3\;+\;u^4H_4\;+\;O(u^5)
\end{equation}
on that shell. This paper is about $H_4$.

\subsection{What was already settled}

Three facts fix the setting and are used throughout. Each is a theorem, with
its own proof recorded elsewhere in this program and its own machine check.

\begin{enumerate}\itemsep3pt
\item \textbf{The shell is isolated, uniformly in the volume.} Within the
  trivial-flux sector the gauge-invariant electric spectrum below
  $\tfrac52\CF$ is exactly $\{0,2\CF\}$ at every $L\ge3$, with
  $\CF=(N^2-1)/(2N)$; the charge-odd part of the $2\CF$ eigenspace is exactly
  the $3L^3$-dimensional one-plaquette shell, and the external margin is at
  least $\CF/2$ (sharpened to $\gamma=\CF$ for $N=3$ and $N\ge5$, and
  $\gamma=5/4$ at $SU(4)$ alone, where the accident $4\cdot\CF/2=2\CF$ occurs).
  The proof is a Casimir census plus a parity lemma: every $5$-cycle on the
  cubic complex wraps, and its flux is winding times $N$-ality, so the
  five-link shelf carries no fundamental flux. The winding energies and the
  three margin identities are symbolic in $N$; the centre-neutral Casimir floor
  and the step that only $F$ and $\bar F$ sit below the window are enumerations,
  machine-verified for $N=3,\dots,7$.\chk{below 5 C\_F/2 the trivial-flux electric spectrum is exactly 0 and 2 C\_F}
\item \textbf{The second-order symbol is the coboundary.} The four shared-link
  fusion channels ($\mathbf1$, $\adj$, $\Sym^2F$, $\Lambda^2F$) carry weights
  $d_\rho/N^2$ --- an order-two Weingarten computation, \emph{not} an isotropy
  hypothesis --- and a link-centre classification shows those four channels
  exhaust every second-order off-diagonal process. Hence
  \begin{equation}\label{eq:t}
    H_2=t_N\,BB\dg,\qquad B=\partial_2\dg,\qquad
    t_N=\frac{2N(N^2-4)}{(N^2-1)(2N^2-1)(4N^2-9)}>0\quad(N\ge3),
  \end{equation}
  with $t_3=5/612$, and the Bloch spectrum of $BB\dg$ is $\{0,q,q\}$ with
  $q(k)=4\sum_{j}\sin^2(k_j/2)$.\chk{t\_2 = 0 and t\_3 = 5/612}
\item \textbf{The flat band is homological.} $\ker\partial_2$ on the periodic
  $L^3$ complex has dimension $L^3+2$; the $\Gamma$ block is $b_2(T^3)$, not a
  coincidence of coefficients, and the flatness is protection by $\partial_2$.
  The boundary of one elementary cube is an exact zero mode of the face Gram
  at every $L$, with $L$-independent squared norm $6$; its normalised Fourier
  sum is an exact carrier at every nonzero allowed momentum.
  Against that, an exact no-go: every $\operatorname{im}\partial_3$ vector has
  exactly zero zone-average, so the whole $k=0$ carrier content is carried by
  the wrapping sheets, and a sheet's normalised overlap with any fixed local
  observable is exactly $1/L$. Every finitely supported exactly-flat operator
  is therefore $\Gamma$-blind.\chk{im d3 is Gamma-blind, and the sheets carry the whole zone-average}
\end{enumerate}

The homological statements are verified at finite $L$ (through $L=5$; open
windows $n=2,3,4$) rather than by induction in $L$; the unbounded-lattice step
that compact $2$-cycles are exactly cube boundaries is classical Alexander
duality, named as an input rather than proved here. The symbolic statements ---
the boundary-factorised rigidity identity, the winding energies and margins,
and the zone minimum --- are uniform in $L$.

Everything above is order $u^2$. At order $u^4$ the carrier acquires
dispersion, and the shape of that dispersion is the subject here.

\subsection{The premises, stated once}\label{sec:premises}

Everything below about the \emph{fourth} order is downstream of four premises.
They are not hidden in the derivation, and each is an open route in this
program's register rather than a theorem.

\begin{hypothesis}[the fourth-order form]\label{hyp:pvp}
When $PVP=0$, the Hermitian fourth-order effective Hamiltonian on the
degenerate manifold is $D-\tfrac12\{H_2,V_2\}$ in the standard notation of
degenerate perturbation theory. This defines the object $H_4$ that the whole
paper is about.
\end{hypothesis}

\begin{hypothesis}[the $E_0$ eigenspace]\label{hyp:e0}
On the reachable words, the $E_0$ eigenspace is exactly the plaquette span, so
its components may be dropped from the resolvent. This was verified
computationally on $224$ dropped components; it is not derived. It is shared by
every implementation below, so a fault here is common-mode and propagates
identically into each.
\end{hypothesis}

\begin{hypothesis}[completeness of the cumulant list]\label{hyp:complete}
The fourth-order connected decomposition is exhausted by the recorded cluster
list. This rests on a charge-counting argument alone. Its named falsifier is
cheap and untried: enumerate the fourth-order support at $N=5$ from the engine
and compare it with the $189$ keys of the Hodge form.
\end{hypothesis}

\begin{hypothesis}[the vacuum intermediates and the rank-generic resolvent]
\label{hyp:resolvent}
The vacuum intermediates are treated as recorded, and the resolvent is taken
rank-generic. At $N=4$ the closed forms additionally rest on the engine's
determinant trick (pseudoinverse Weingarten at $n>N$ with an $\varepsilon$
insertion), a path no independent lens has audited.
\end{hypothesis}

One convention belongs here too, because every cumulant obeys it: the effective
Hamiltonian is \emph{not} cluster-additive. $H_{\mathrm{eff}}-eI$ is. Möbius
subtraction therefore acts on the vacuum-subtracted operator throughout, and
that is what makes the connected cumulants well defined.

\subsection{The shape ansatz, and the two coefficients that vanish}

Write $a_j=4\sin^2(k_j/2)$, so $q=a_1+a_2+a_3$, and let $e_1=q$, $e_2$, $e_3$
be the elementary symmetric functions of $(a_1,a_2,a_3)$. The carrier direction
is $d$ with $d\dg d=q$, so the fourth-order band correction is
$\varepsilon_4=(d\dg H_4 d)/q$, which is where the $1/q$ in the ansatz comes
from. Cleared of that denominator, the fourth-order shape is the linear form
\begin{equation}\label{eq:ansatz}
  q\,\varepsilon_4 \;=\; c_0\,q \;+\; A\,q^2 \;+\; 4C\,e_2
                  \;+\; B\,(q\,e_2) \;+\; D\,e_3 ,
\end{equation}
with degrees $1,2,2,3,3$ in the $a_j$. Since $a_j\sim L^{-2}$ near $\Gamma$,
the degree-$3$ monomials $q\,e_2$ and $e_3$ constitute the $L^{-4}$ tier. The
recorded kernels have
\begin{equation}\label{eq:collapse}
  A=\tfrac{5}{48},\qquad B=0,\qquad D=0 ,
\end{equation}
and the vanishing of the pair $(B,D)$ is what this program calls the
\emph{tier collapse}. It is not kinematic. The two-hop enumeration does
generate the degree-$3$ monomials --- the rank-five numerator span contains
both --- so the dynamics sets their coefficients to zero.

\begin{remark}[a retraction that shapes this paper]\label{rem:retraction}
An earlier route derived the collapse from a vertex-counting degree bound:
$2r$ vertices at order $r$ would bound the numerator degree by $r$, forbidding
$B$ and $D$ at fourth order and predicting the $L^{-4}$ tier first at sixth
order. That bound is false and was withdrawn. The numerator is $d\dg H_4d$,
not $H_4$; the carrier projection contributes one further $d$ and one further
$\bar d$, so four vertices bound the \emph{entries} of $H_4$ at $a$-degree $2$
but bound the \emph{numerator} at $3$ --- exactly the degree the count was
meant to exclude. An explicit witness is
$d\dg\operatorname{diag}(a_i^2)d=\sum a_i^3=q^3-3qe_2+3e_3$, nonzero in $e_3$.
The failed route is kept on the record because it tells the next attempt where
the trap is: nothing below is a vertex count. Every statement in
Section~\ref{sec:hodge} is a Laurent-polynomial identity of named operators.
\end{remark}

\subsection{Results}

\paragraph{(A) A Hodge decomposition of the fourth-order kernel
(Section~\ref{sec:hodge}).}
Let $\Ldn=\partial_2\dg\partial_2$ and $\Lup=\partial_3\partial_3\dg$ be the
down and up plaquette Laplacians, and let $\Ssq=\Ldn-4I$ be the signed
shared-link adjacency. Then, with $\tilde\nu=\nu+4u$, $\tilde\pi=\pi+2u$,
$\tilde\sigma=\sigma-12u$ and $R$ the cross-plane half of $\Ssq$,
\[
  H_4 \;=\; -\tilde\nu(\Lup-2)\;+\;u\,\Ssq^2\;-\;\tilde\pi\,\Ssq\;+\;\tilde\sigma I
        \;-\;2C_{\mathrm{shp}}\,R ,\qquad \tilde\nu=-\tfrac{5}{48},
\]
exactly on all $189$ records of the historical kernel (Theorem~\ref{thm:hodge}).
The $144$ records the two rival kernels agree on are $u\,\Ssq^2$ exactly: the
two-hop sector enters as the \emph{square} of the second-order operator.

\paragraph{(B) The tier collapse as an algebra statement
(Section~\ref{sec:collapse}).}
At each Bloch point with $q>0$ the plaquette fibre splits as
$\C\psi\oplus\psi^\perp=\ker\Ldn\oplus\ker\Lup$, because $\Ldn+\Lup=qI$ per
plane component. So the Feshbach complement is exactly the projector onto
$\ker\Lup$, and $Q\Ldn\psi=Q\Lup\psi=0$: neither Hodge generator leaves the
carrier, and every excursion costs one insertion of $R$
(Proposition~\ref{prop:feshbach}). Computing carrier symbols word by word,
the actual fourth-order support $\{I,\Ssq,\Ssq^2,\Lup,R\}$ has symbols in
$\operatorname{span}\{q,q^2,e_2\}$, hence $B=D=0$ (Theorem~\ref{thm:support}).
Two sharper facts come with it, and both are recorded as falsifiers rather than
as consequences: $R$-degree alone does \emph{not} give the collapse, because
$\sigma(UR)=-2q\,e_2$ populates $B$ with only one $R$
(Finding~\ref{find:ur}); and two insertions unlock the $L^{-4}$ tier at a
locked ratio, $\sigma(RR)=q\,e_2+3e_3$, so $D=3B$ in the $R^2$ channel
(Proposition~\ref{prop:rr}).

\paragraph{(C) The universality of the two-hop weight
(Section~\ref{sec:univ}).}
The decomposition (A) has exactly one dynamical input: that one weight $u$
multiplies every two-hop chain geometry, so that $u_2=2u$ and $D=-6u+3u_2=0$
for \emph{any} value of $u$. That input is now reduced to a finite identity of
Haar integrals. Under a path-reduction lemma a private path of $k$ links is one
Haar link carrying $k$ times the single-link $H_0$, so a cluster reduces to two
shared links and a few weighted private ones; on the reduced clusters the
straight and the L-shaped chain differ by an explicit map $\varphi$ on words
(delete $B,B\dg$, rename $A\mapsto C$, reverse the second shared link), and
\[
  \int w \;=\; \int\varphi(w)
\]
for every word $w$ in the family the Fierz rewirings generate
(Lemma~\ref{lem:univ}). The lemma is verified by exhaustion at the
multiplicities fourth order reaches, with two controls that locate its scope: on
arbitrary rewirings of the same letters $\varphi$ does \emph{not} preserve the
integral, and cutting the letters no history can bring together breaks it too.

\paragraph{(D) Closed forms at every rank, and $\beta_N$ from three cumulants
(Section~\ref{sec:qn}).}
A Wilson-loop calculus written from scratch --- $H_0$ by the Fierz identity as
a rewiring of index ports, Haar integrals by the $U(N)$ Weingarten function as
the Gram pseudoinverse on $S_n$, the resolvent by per-link irrep projectors ---
is run over the rational function field $\Q(N)$ rather than rank by rank. Every
fourth-order cumulant then comes out as a \emph{single rational function of
$N$}, derived rather than reconstructed: the two-hop weight, the single-contact
and fan dressings of both pairs, the corner, and both cube completions
(Theorem~\ref{thm:closed}). The coplanar dressings are the exact negatives of
the perpendicular ones in the C-odd sector, the pair clusters cancel, and the
eleven-cumulant assembly collapses to
\begin{equation}\label{eq:betaassembly}
  \beta_N=-16u+32\,d_{\mathrm{single}}-16\,d_{\mathrm{corner}}
          +\frac{848}{N(N^2-1)^3}
        =\frac{P_{17}(N^2)}{N\,R_{20}(N^2)} ,
\end{equation}
a polynomial identity of degree $34$ in $N$ (Theorem~\ref{thm:beta}), which is
proof-checked in Lean~4 with no \texttt{sorry} and standard axioms only. This
replaces a previous conditional identification: the same equality had been
established at $N=4,\dots,70$ by rank sweep and could be stated for all $N$
only under a degree hypothesis. The hypothesis is now the computed degree.

\paragraph{(E) The off-axis coefficient (Section~\ref{sec:c2}).}
The disagreement between the two recorded fourth-order kernels ---
$C_{\mathrm{old}}=-0.0480863\ldots$ against $C_{\mathrm{new}}=-0.0202133\ldots$
--- was the program's last open contradiction. Three steps close it.
\emph{(i)} Clearing the $1/q$ makes \eqref{eq:ansatz} a linear identity between
Laurent polynomials in $z_j=e^{ik_j}$, so the coefficients are \emph{solved
over the whole zone} instead of fitted at four points: $A=5/48$, $B=D=0$,
residual identically zero (Theorem~\ref{thm:solved}). \emph{(ii)} In that basis
the $189$ records carry only six weight magnitudes, one per cubic orbit, and
the dispute is the single scalar
$C_{\mathrm{shp}}=-5/96-u-(\rho+\tilde\pi)/2$. \emph{(iii)} A third,
independent implementation of the corner cluster agrees with the assembly to
the digit, and the original pipeline's own word ledger --- recovered from the
corpus and re-run --- agrees on \emph{every} cluster of the rotation record
except the adjacent-face cube completion, where it holds $8$ of the $24$
insertion orderings, each with exactly the weight the multi-loop law assigns
it. The sixteen it lacks sum to $-25/512$, so in the kernel's basis
\begin{equation}\label{eq:cshp}
  C_{\mathrm{shp}}=C_{\mathrm{historical}}+\frac{25}{1024}
  =-\frac{13035490122347}{550663802582400}=-0.0236723\ldots
\end{equation}
(Theorem~\ref{thm:c2}). Both recorded kernels are wrong: the historical one by
those sixteen orderings, the other on $u$, $\pi$ and $\rho$ alike.

\paragraph{(E$'$) Why the value matters: one angular invariant
(Section~\ref{sec:aniso}).}
The resolved coefficient is not bookkeeping. Let $x_i=a_i/q$ and
\[
  V(x)\;=\;\sum_i x_i^3-\Bigl(\sum_i x_i^2\Bigr)^{2}
        \;=\;\sum_{i<j}x_ix_j\,(x_i-x_j)^2 .
\]
Then the \emph{entire} fourth-order coupling of the carrier to its two
transverse bands is one nonnegative angular invariant,
\[
  Q H p=-2C_{\mathrm{shp}}\,QDp ,\qquad
  \bigl\|QHp\bigr\|^2 \;=\; 4C_{\mathrm{shp}}^2\,q^2\,V(x)
\]
exactly --- an $85$-coefficient cleared Laurent identity, with the shape
identity $V=e_2+3e_3-4e_2^2$ on the simplex proof-checked in Lean
(Theorem~\ref{thm:aniso}). Two consequences follow. The nodal set of the
mixing is exactly the isotropic directions --- axes, equal face diagonals,
equal body diagonals (Proposition~\ref{prop:nodal}) --- and the induced
sixth-order coefficient is a definite negative number with a locked shape
ratio,
\[
  -\frac{4C_{\mathrm{shp}}^2}{t_3}
  =-\frac{169924002729806205028788409}{619343593697933825385600000}
  =-0.2743614\ldots,\qquad \text{shape } 1:3:-4 ,
\]
(Proposition~\ref{prop:sixth}). Because the coefficient enters as
$C_{\mathrm{shp}}^2$, the two recorded branches differ here by a factor of
about $4.1$: deciding \eqref{eq:cshp} decides a sixth-order number.

\paragraph{(F) What follows for the $SU(3)$ shape row
(Section~\ref{sec:su3}).}
$\alpha=4A=5/12$ and the axial coefficient $\lambda_X$ are unchanged; the
off-axis quantities move by exactly the crosswalk multiples of $25/1024$:
$\beta_3$, $W_4$ and $\lambda_R$ by $25/64$ and $\lambda_M$ by $25/128$
(Corollary~\ref{cor:table}). In particular $\beta_N$ at $N=3$ is now
$8A+16C_{\mathrm{shp}}$ exactly, so nothing distinguishes $N=3$ from $N\ge4$ in
the fourth-order shape coefficient --- the ``$SU(3)$ exception'' recorded in
the corpus was the cube shortfall.

\subsection{What this paper does not claim}\label{sec:notclaim}

Stated once here and in detail in Section~\ref{sec:scope}. Every result above
is a statement about an effective operator obtained by degenerate perturbation
theory on a projected finite-dimensional shell, at \emph{fixed lattice
spacing}. It is not a statement about the continuum theory, and it is not a
mass gap.

Two of the obstructions are theorems rather than cautions. In the Hamiltonian
coordinate $u=g_H^{-4}$, an asymptotically free trajectory has $u(a)\to\infty$
while the proved small-$u$ persistence island requires $u<u_*$; the two
hypothesis sets are eventually \emph{disjoint}, so no tail of a continuum
trajectory lies in the proved domain, and evaluating the existential constants
would not change that. And no nonzero finite strong-coupling Taylor truncation
yields a finite continuum mass at all, since $g_H^2P(u(a))/a$ diverges for
every nonzero polynomial $P$.

\section{Setting and notation}\label{sec:setting}

\subsection{The complex, the coboundary, and the two Laplacians}

Let $C_1,C_2,C_3$ be the oriented cellular chain groups of $\Lambda$ over
$\Q$ --- links, plaquettes, cubes --- with boundary maps
$\partial_2:C_2\to C_1$ and $\partial_3:C_3\to C_2$, and $\partial_2\partial_3=0$.
Write $B=\partial_2\dg$ for the coboundary on plaquettes and
\begin{equation}\label{eq:laplacians}
  \Ldn=\partial_2\dg\partial_2=B\dg{}\!^{\dagger}\!B\dg,\qquad
  \Lup=\partial_3\partial_3\dg,\qquad \Ldn\Lup=\Lup\Ldn=0 .
\end{equation}
The Hodge decomposition of $C_2$ is
$\operatorname{im}\partial_2\dg\oplus\ker\Delta_2\oplus\operatorname{im}\partial_3$
with $\Delta_2=\Ldn+\Lup$. On the Bloch fibre at momentum $k$ each of these is
$3$-dimensional and
\begin{equation}\label{eq:fibre}
  \Ldn+\Lup=q(k)\,I_3 ,\qquad q(k)=4\sum_{j=1}^{3}\sin^2(k_j/2)=\sum_j a_j .
\end{equation}
This last identity is the one that makes Section~\ref{sec:collapse} work: the
two Laplacians are complementary spectral projections scaled by $q$, so on each
fibre with $q>0$ they commute with, and are diagonal in, the same
two-dimensional splitting.

The signed shared-link adjacency is $\Ssq=\Ldn-4I$; the carrier is the
one-dimensional $\ker\Ldn$ in each fibre with $q>0$, spanned by $\psi=Bd/\|\cdot\|$
where $d$ is the Fourier transform of a radius-one cube boundary, and
$\Ssq\psi=-4\psi$.

\subsection{Records, orbits, and the six amplitudes}\label{sec:orbits}

The fourth-order kernel is recorded as $189$ matrix elements between plaquette
pairs, labelled by their displacement and relative orientation. Grouping them
by displacement gives the ``blocks'' of earlier accounts; grouping them by
\emph{weight magnitude} gives a smaller and much more informative basis.

\begin{proposition}[six orbits]\label{prop:orbits}
The $189$ records carry exactly six distinct weight magnitudes, one per cubic
orbit: the $132$-record skeleton at $u$, the $12$-record same-plane $(0,1,1)$
in-plane orbit at $u_2$, the $24$-record cross-plane rotation orbit at $\rho$,
the $12$-record coplanar orbit at $\pi$, the normal orbit at $\nu$, and the
on-site amplitude $\sigma$.\chk{the 189 records carry exactly six weight magnitudes}
\end{proposition}

Each orbit's carrier projection is closed form in $e_1,e_2,e_3$, and each is an
exact identity on that orbit's own records:
\begin{equation}\label{eq:orbitforms}
\begin{aligned}
  \sigma&\longmapsto \sigma\,e_1, &\quad
  \nu&\longmapsto \nu\,(2e_1-e_1^2+2e_2), &\quad
  \pi&\longmapsto \pi\,(4e_1-2e_2),\\
  u_2&\longmapsto u_2\,(4e_1-4e_2+3e_3), &\quad
  \rho&\longmapsto -2\rho\,e_2, &\quad
  u&\longmapsto u\,(12e_1-4e_1^2+12e_2-6e_3).
\end{aligned}
\end{equation}
Five of the six follow from the orbit's displacement set in two lines. The
rotation one follows because
$S_{PQ}=-\varepsilon_P\varepsilon_Q\,\rho\,\bar d_{n(Q)}\overline{\bar d_{n(P)}}$
exactly, so the carrier kills the phases and leaves
$-\rho\sum_{m\ne n}s_ms_n$.\chk{every orbit's carrier projection is closed form in e1, e2, e3}

Two consequences are immediate from \eqref{eq:orbitforms} and are worth
separating, because they were previously read as one phenomenon.

\begin{corollary}[$B$ is unpopulated; $D$ is one equality]\label{cor:bd}
Not one of the six forms in \eqref{eq:orbitforms} contains the monomial
$e_1e_2$. Since $B$ is the coefficient of $e_1e_2$, the coefficient $B$ is
\emph{unpopulated} rather than cancelled --- there is nothing to cancel. And
only two orbits carry $e_3$: the skeleton at $-6u$ and the same-plane $(0,1,1)$
in-plane orbit at $+3u_2$. Hence
\[
  D=-6u+3u_2 ,
\]
which vanishes identically once $u_2=2u$, for \emph{any} value of $u$. It is
not a numerical cancellation between sectors.\chk{the tier collapse is one identity: D = -6u + 3*u2 with u2 = 2u}
\end{corollary}

This supersedes an earlier displacement-block reading in which $B=0$ appeared
as an integer cancellation $(+1,+1,-2)$ across three blocks and $D=0$ as
$(-3,+6,-3)$ across three others, coupling the translation sector to the
$120$-record rotation sector. Those cancellations are real in that basis; they
are artefacts of the basis. In the orbit basis there is one equality left,
$u_2=2u$, and Sections~\ref{sec:collapse}--\ref{sec:univ} are about it.

\begin{remark}[why $u_2=2u$ is the natural statement]
Write $X=2\cos k_i$, $Y=2\cos k_j$, $Z=2\cos k_n$ for a plane with own axes
$\{i,j\}$ and normal $n$. Everything the plane sends to itself beyond nearest
neighbour is sixteen records: $(0,0,2)$ at $u$ giving $(X^2-2)+(Y^2-2)$, mixed
$(0,1,1)$ at $u$ giving $(X+Y)Z$, and in-plane $(0,1,1)$ at $u_2$ giving
$(u_2/u)XY$. Their sum is $u[(X+Y)(X+Y+Z)-4]$ \emph{if and only if} the cross
term is $2XY$. So $u_2=2u$ is exactly the condition that the same-plane
long-range symbol \emph{factorises} rather than being an inhomogeneous
quadratic --- and with the product form $D=0$ holds separately in the
same-plane sector and in the $96$-record cross-plane
skeleton.\chk{u2 = 2u is the cross term of a perfect square}
\end{remark}

\section{The Hodge decomposition of the fourth-order kernel}\label{sec:hodge}

The factorisation in the last remark is a square, and the square is the second-order
operator's.

\begin{theorem}[the kernel is Hodge plus one operator]\label{thm:hodge}
With $\tilde\nu=\nu+4u$, $\tilde\pi=\pi+2u$, $\tilde\sigma=\sigma-12u$, and $R$
the cross-plane half of $\Ssq$ (equivalently, minus the cross-plane half of
$\Lup$),
\begin{equation}\label{eq:hodge}
  H_4=-\tilde\nu\,(\Lup-2)\;+\;u\,\Ssq^2\;-\;\tilde\pi\,\Ssq\;+\;\tilde\sigma\,I
      \;-\;2C_{\mathrm{shp}}\,R
\end{equation}
holds on all $189$ records of the historical kernel exactly, with
$\tilde\nu=-5/48$; and on the second recorded kernel to $10^{-13}$ relative,
with the same $\tilde\nu$.\chk{H4 = -nu~(L\_up - 2) + u S\_sq\^{}2 - pi~ S\_sq + sigma~ - 2 C\_shp R exactly: C\_shp is the coefficient of the one non-Hodge operator}
\end{theorem}

The proof is incidence algebra on records. Three of its ingredients deserve to
be stated separately because they carry the content.

\begin{proposition}[the disputed orbits are $\Ssq$]\label{prop:signs}
The unit sign pattern of the two disputed orbits --- the $24$ cross-plane
records of $\rho$ and the $12$ coplanar records of $\pi$ --- is exactly the
off-diagonal of $\Ldn$, i.e.\ of $\Ssq=\Ldn-4I$.\chk{no plane-basis convention flips rho or pi: only +-1 keeps the 144 agreed records}
\end{proposition}

\begin{proposition}[the agreed records are the square]\label{prop:square}
As an operator product on records, $\Ssq^2$ is the $132$-record skeleton plus
the $12$-record doubled orbit at unit weight --- so $u_2=2u$ is precisely the
statement that the doubled orbit counts the two coplanar paths --- plus a
diagonal shadow: $-4$ on the normal keys, $-2$ on the in-plane keys, $+12$ on
site. The $144$ records the two rival kernels agree on in shape are $u\,\Ssq^2$
exactly, and the $-4u$ the normal orbit carries beyond $-5/48$ is that shadow,
not a correction to the cube channel.\chk{the 144 agreed records are u S\_sq\^{}2, the shared-link adjacency squared}
\end{proposition}

\begin{corollary}[the collapse on the carrier]\label{cor:carrier}
On the carrier, $\Ldn\psi=0$ and $\Lup\psi=e_1\psi$, so in \eqref{eq:hodge}
every term but the last acts as a scalar. Hence $A=-\tilde\nu=5/48$ with no $u$
in it, $B=D=0$ with nothing to cancel, and the entire $e_2$ content is the
projection of $-2C_{\mathrm{shp}}R$ alone, which is $4C_{\mathrm{shp}}e_2$. In
particular the carrier is an exact eigenvector of $H_4$ if and only if
$C_{\mathrm{shp}}=0$.\chk{C\_shp = -5/96 - u - (rho + pi)/2, exactly}
\end{corollary}

Theorem~\ref{thm:hodge} answers the ``why a product'' question of
Section~\ref{sec:orbits}: it is a square. And it isolates the whole of the
fourth-order dispute in one scalar, the weight of the single operator outside
the algebra generated by the two Laplacians.

\begin{remark}[the single weight on $\Lup$ is a reparametrisation]
Equation~\eqref{eq:hodge} writes the cube channel as $-\tilde\nu(\Lup-2)$, one
weight for every pair of faces of a cube. Computed from primitives, the
cube-completion channel --- the four-insertion histories that are the cube's
other four faces --- is $-5/48$ for \emph{opposite} faces, reproducing
$\tilde\nu$ exactly, and $53/768$ for \emph{adjacent} faces. So a single weight
on $\Lup$ is not a mechanism; on the historical records it is a
reparametrisation that absorbs the difference into the rotation orbit. Any
proposed mechanism for the collapse has to produce $-160$ and $-106$ together
(Section~\ref{sec:qn}, \eqref{eq:cubes}).\chk{the cube completion from the third engine: -5/48 between opposite faces, -53/768 between adjacent ones}
\end{remark}

\section{The tier collapse is a statement about the algebra}\label{sec:collapse}

Corollary~\ref{cor:carrier} derives $B=D=0$ from the \emph{form}
\eqref{eq:hodge}. This section derives it from the \emph{support}: which words
in the generators can appear at all, and what each contributes to the carrier
symbol. That is stronger, because it does not presuppose that the assembled
coefficients take the values they take.

\subsection{The Feshbach complement is $\ker\Lup$}

\begin{proposition}[the complement, and the cost of leaving]\label{prop:feshbach}
At each Bloch point with $q>0$ the three-dimensional plaquette fibre splits as
\[
  \C\psi\;\oplus\;\psi^{\perp}\;=\;\ker\Ldn\;\oplus\;\ker\Lup ,
\]
because $\Ldn+\Lup=qI$ per plane component and $\Ldn\psi=0$. Consequently the
Feshbach projector $Q=1-P_\psi$ is exactly the projector onto $\ker\Lup$, and
\[
  Q\,\Ldn\,\psi=Q\,\Lup\,\psi=0 .
\]
Neither Hodge generator can leave the retained sector: every excursion off the
carrier costs one insertion of the one non-Hodge operator $R$, whose amplitude
in \eqref{eq:hodge} is $-2C_{\mathrm{shp}}$. Moreover the $R$-excitation is
up-harmonic, $\Lup(QR\psi)=0$ exactly.\chk{the R-excitation is up-harmonic: L\_up (Q R psi) = 0 exactly}
\end{proposition}

The identities behind Proposition~\ref{prop:feshbach} are
$\Ldn+\Lup=qI$, $\Ldn\Lup=0$, $\Lup=\psi\psi\dg$ and $\psi\dg\psi=q$, each
verified on all $21$ plane-diagonal records; the $R$-excitation
$\phi=qR\psi-\langle\psi,R\psi\rangle\psi$ is nonzero, carrier-orthogonal and
satisfies $\Lup\phi=0$.

\begin{proposition}[when a defect can appear at all]\label{prop:defect}
Of the $39$ words of length at most three in $(\Ssq,U,R)$ --- writing $U$ for
$\Lup$ --- exactly seven have a nonzero Feshbach defect:
$RR$, $RRR$, $RRS$, $RRU$, $RSR$, $SRR$, $URR$. Every one of them contains two
$R$'s with no $U$ between them; $RUR$ factorises, because $\Lup\phi=0$, and
$RSR$ does not. The rule is about the intermediate excursion, not about
counting $R$'s. Moreover, every word of length at most three with at most one
$R$ obeys the closed formula
\begin{equation}\label{eq:wordformula}
  \sigma(W)=(-4)^{\#S}\,(-2)^{\#R}\,q^{\,\#U+1-\#R}\,e_2^{\,\#R},
\end{equation}
verified as cleared Laurent identities.\chk{words of length at most three with at most one R obey the exact carrier formula}
\end{proposition}

Neither statement is an all-length theorem; both are scoped to lengths one
through three, which is what fourth order reaches.

\subsection{The collapse}

The algebraic core of this section is now proof-checked: a Lean module
formalises the Hodge relations $\Ldn+\Lup=q\,\mathrm{id}$ and
$\Ldn\Lup=0$ with their consequences $\Ldn^2=q\Ldn$, $\Lup\Ldn=0$ and
$\Lup^2=q\Lup$; the homological protection $\Ldn\psi=0$, $\Lup\psi=q\psi$,
$\Ssq\psi=-4\psi$ and $Q\Ldn\psi=Q\Lup\psi=Q\Ssq\psi=0$; the up-harmonicity
and carrier-orthogonality of $\phi=QR\psi$; the cleared defect table for all
$39$ words of length at most three; and the selection rule of
Proposition~\ref{prop:defect}.\chk{for q > 0, Q projects onto ker L\_up; the cleared Hodge identities hold everywhere}

\begin{theorem}[the actual support forces $B=D=0$]\label{thm:support}
Write $\sigma(w)$ for the carrier symbol of a word $w$ in the generators. For
the actual fourth-order support $\{I,\Ssq,\Ssq^2,\Lup,R\}$,
\[
  \sigma(w)\in\operatorname{span}_{\Q}\{q,\;q^2,\;e_2\}\qquad\text{for every }w,
\]
so no word populates $e_1e_2$ or $e_3$, and therefore $B=D=0$. The statement
requires only $\Ldn\psi=0$ and $\Lup\psi=\lambda\psi$, and both hold in three
geometries.\chk{the actual fourth-order support I, U, S, S\^{}2, R forces B = D = 0}
\end{theorem}

This is the mechanism, at the level of the operator algebra. It is not a
statement about every $R$-linear kernel, and the next two facts say why the
weaker readings fail --- both are recorded as findings, i.e.\ as checks that
\emph{assert} a discrepancy rather than hide it.

\begin{finding}[$R$-degree alone is not enough]\label{find:ur}
$\sigma(UR)=\sigma(RU)=-2q\,e_2$. A single $R$ therefore populates the
$B$-monomial. The selection rule ``$B$ and $D$ are excluded by $R$-degree''
is false as stated; the correct statement is
Theorem~\ref{thm:support}, about the actual support.\chk{FINDING: one R does not exclude B; sigma(UR) = sigma(RU) = -2 q e\_2}
\end{finding}

\begin{proposition}[two $R$'s unlock the tier at a fixed ratio]\label{prop:rr}
$\sigma(RR)=q\,e_2+3e_3$ exactly. So the $R^2$ channel carries both top-degree
monomials, locked at $B:D=1:3$; i.e.\ $D=3B$ in that channel. The $L^{-4}$
tier is not forbidden --- it is unpopulated at fourth order and enters as soon
as two $R$'s can meet.\chk{two R insertions unlock the L\^{}-4 tier locked at B : D = 1 : 3}
\end{proposition}

\begin{finding}[a word outside the shape span]\label{find:rur}
$\sigma(RUR)=4e_2^2$, whose shape symbol $4e_2^2/q$ lies \emph{outside} the
five-element span of the ansatz \eqref{eq:ansatz}. The corpus's obstruction
hypothesis --- that the obstruction space is spanned by elementary symmetric
polynomials and nothing else appears --- is therefore violated by a word of the
very algebra that produced the fourth-order kernel. Whether the sixth-order
dynamics populates that word is open; this statement is an identity of the
algebra, not a claim about the dynamics.\chk{FINDING: the algebra contains a shape monomial the four-shape ansatz cannot hold}
\end{finding}

Findings~\ref{find:ur} and~\ref{find:rur} are the reason
Remark~\ref{rem:retraction} is in the introduction. A selection rule that
``looks like'' the collapse is not the collapse; the program has already
withdrawn one such rule.

\begin{remark}[the hypothesis Theorem~\ref{thm:support} actually needs]
A ``link regularity'' condition --- that $\Ldn+\Lup$ be scalar --- was proposed
as the requirement. It is not one. The weaker hypothesis, that $\ker\Ldn$ is
one-dimensional and spanned by $\psi$ with $\Lup\psi=\lambda\psi$, already makes
every word in the two Hodge Laplacians act as a scalar on the carrier, and it
holds where the strong condition fails: on the tetrahedron $\lambda=4$ and
$\Ldn+\Lup=4I$ is scalar, while on the pentagonal prism $\lambda=7$ and
$\Ldn+\Lup=\operatorname{diag}(6,6,5,5,5,5,5)$ is not.\chk{the Hodge-word statement needs only L\_down psi = 0 and L\_up psi = lambda psi, and both hold in three geometries}
\end{remark}

\section{The universality of the two-hop weight}\label{sec:univ}

Corollary~\ref{cor:bd} left one equality, $u_2=2u$; Theorem~\ref{thm:hodge}
left one dynamical input, that the two-hop sector enters as $u\,\Ssq^2$ with a
\emph{single} weight. They are the same statement: every two-hop chain
geometry --- coplanar--coplanar, coplanar--perpendicular,
perpendicular--perpendicular --- carries the same fourth-order weight, up to
the incidence sign $S_{PQ}S_{QR}$.

\subsection{Computed first, at three levels}

The statement was established computationally before it was explained, and the
sequence matters because each step narrowed what a proof had to do.

\begin{enumerate}\itemsep3pt
\item \textbf{Numerically, across geometries.} An independent fourth-order
  assembly returns $+u$ on the coplanar chain and $-u$ on the bent chain,
  exactly, with $u=360421351/40327601932800$ at $N=3$; a second, independently
  written implementation (engine $PVP$ assembly, Krylov resolvent) returns the
  same rational on three chain types including the coplanar--perpendicular one,
  where the two hops carry opposite signs and the cumulant carries their
  product. A previously recorded value $4.13\times u$ in one of the two kernels
  is thereby shown to be wrong.\chk{FINDING: the chain amplitude is u = X\_QUANTUM exactly, on the coplanar and the bent chain; the cold kernel's 4.13 u is wrong}
\item \textbf{Channel by channel.} Decomposing the two-hop cumulant by the
  irreps of its intermediate states gives $74$ channels, each a product of the
  resolvent factors its irreps bring ($\Lambda^2$: $2N-3$, $3N-4$; $\Sym^2$:
  $2N+3$, $3N+4$; adjoint: $2N^2-1$, $3N^2-2$; mixed pairs: $2N^2-3$ and
  $4N^2\mp N-4$), summing to $u(N)$ identically; and the coplanar, bent and
  in-plane L chains agree in \emph{every one} of the $74$ channels, up to the
  incidence sign, at every rank $N=3,\dots,30$. Over $\Q(N)$ the same
  statement holds as an identity of rational functions, with the finer
  per-link labelling distinguishing $122$ channels, the extra $48$ vanishing in
  the C-odd sector.\chk{u is universal channel by channel as an identity in Q(N): bent = -coplanar and L = coplanar in every C-odd channel and both equal it in every C-even channel, so u2 = 2u and the tier collapse D = -6u + 3u2 = 0 hold at every rank}
\item \textbf{History by history.} Under the path-reduction lemma below, the
  agreement holds separately for each time-ordered insertion sequence, and
  integrand by integrand.\chk{universality of the two-hop weight holds history by history: the straight and the L-shaped chain agree in every (insertion sequence, channel) term of the direct term, and integrand by integrand within each term}
\end{enumerate}

\subsection{Path reduction}

\begin{lemma}[private paths are single links]\label{lem:reduce}
A private path of $k$ links --- a path none of whose links is shared between the
two plaquettes of interest --- is equivalent, inside every cumulant, to one
Haar link carrying $k$ times the single-link $H_0$. Consequently every
fourth-order cumulant is computed on a \emph{reduced cluster}: two shared links
and a few weighted private ones.\chk{a private-link path is one effective link with k times the single-link H0: every three-cluster cumulant of the beta\_N assembly is the same, channel by channel in both sectors, on the reduced cluster as on the full one}
\end{lemma}

Reduction is what makes the universality statement finite. On the reduced
clusters the two chains share $P=\Tr(U_1W_P)$ and $R=\Tr(U_2W_R)$; the straight
middle face is $\Tr(U_1\dg AU_2B)$, with two private links of weight one, and
the L middle face is $\Tr(U_1\dg CU_2\dg)$, with one private link of weight two
in $A$'s place and the second shared link traversed the other way.

\subsection{The lemma}

\begin{definition}[the map $\varphi$]
On words, $\varphi$ deletes $B$ and $B\dg$, joining predecessor to successor;
renames $A,A\dg$ to $C,C\dg$; and reverses the second shared link.
\end{definition}

\begin{lemma}[universality]\label{lem:univ}
Let $w$ be any word obtained from $P^a\bar P^a R^b\bar R^b X\bar X$ by
permuting the out-targets of the letters on one shared link among themselves
(the Fierz swaps act on like pairs; every permutation is allowed here) and then
applying Fierz cuts to unlike pairs of shared-link letters that a history can
bring into one resolvent-projected state --- so never $X$'s own letter against
$\bar X$'s --- each cut rewiring the two wires across and counting an emptied
cycle as a loop of $N$. Then
\[
  \int w \;=\; \int \varphi(w).
\]
\end{lemma}

\begin{proof}[Status of the proof, stated exactly]
The lemma is \emph{verified}, not proved, and the verification is exhaustive
only in part. Exhaustive: every swap word for $a=b=1$ ($576$ words, at two
ranks and over $\Q(N)$) and for $(2,1)$ and $(1,2)$ ($17{,}280$ each); every
cut word for $a=b=1$ ($384$). Sampled: the swap family of $(2,2)$ at
$5{,}129$ of $518{,}400$ words (about one percent), and the reachable
swap-then-cut families of $(2,1)$ and $(1,2)$ at $2{,}052$ words each. Every
word tested passes. Since $(2,2)$ is a multiplicity fourth order reaches, the
lemma is not exhausted even at fourth order --- this is the gap between
``checked'' and ``proved'' at the order this paper is about, and closing it
does not require new theory, only the remaining $99\%$ of one family.

Where the verification is exhaustive it is proof for each word, because the
integrals are rational functions of $N$ of degree bounded by the Weingarten
sums involved, so agreement at two ranks determines them. Two controls locate
the scope: on arbitrary wirings of the same letters $\varphi$ does \emph{not}
preserve the integral, and cutting $X$'s letter against $\bar X$'s --- which no
history can do, since they sit on opposite sides of the amplitude --- breaks it
as well. The identity is a property of the family the Fierz rewirings generate,
not of the map.\chk{the universality lemma, by exhaustion on the fourth-order family: phi (delete the second private letter of the straight middle face, keep the first as the weight-two link) maps every history's straight integrand onto its L integrand with the same coefficient, and preserves the Haar integral on every word the Fierz swaps and cuts can generate from the face words}
\end{proof}

\begin{remark}[what is a fact, and what is the explanation of it]\label{rem:factvsmech}
The universality of $u$ is \emph{established} independently of
Lemma~\ref{lem:univ}, as an exact identity of rational functions over $\Q(N)$:
the coplanar, bent and in-plane L chains agree in every one of the channels,
up to the incidence sign, as elements of $\Q(N)$. That is the fact, and it is
what $u_2=2u$ and $D=0$ rest on. Lemma~\ref{lem:univ} is the \emph{mechanism} ---
it says why, at the level of Haar integrals of individual words --- and it is
the part that is sampled at $(2,2)$. Nothing in Section~\ref{sec:qn} or
Section~\ref{sec:c2} depends on the lemma; only the claim to have explained the
collapse does.
\end{remark}

\begin{remark}[two recorded near-misses]
The route to Lemma~\ref{lem:univ} produced two negative results worth keeping.
First, the identity is \emph{not} local to the middle face: in $14$ of the
$124$ terms it fails term-locally and closes only once the end faces' hairpins
close the shared-link letters. Second, an earlier staging of the correspondence
matched only $52$ of $128$ cases; the map was right and the staging was wrong,
and the two-hop weight's own record required the time-reversed reading of the
channel labels before the correspondence became a bijection.
\end{remark}

\begin{theorem}[the collapse, assembled]\label{thm:mechanism}
In every surviving fourth-order history the middle face appears once on each
side of the amplitude: a history with $X\bar X$ on one side leaves the private
links in the adjoint after the resolvent, and the final Haar integral kills it.
So the private letters are never projected, the Fierz rewirings act on the
shared-link letters only and identically in both geometries, and every
integrand of the straight chain is a word of the family of
Lemma~\ref{lem:univ}. In all $124$ $(\text{sequence},\text{channel})$ terms,
$\varphi$ maps the $536$ straight integrands bijectively onto the L integrands
with the same coefficient. Hence every term of the two-hop weight agrees
between the geometries, hence every channel, hence $u_L=u_{\mathrm{straight}}$,
hence $u_2=2u$, hence $D=0$; and $B=0$ is the carrier projection of
Corollary~\ref{cor:bd}. The tier collapse at fourth order rests on the
incidence algebra of the cubical complex and a finite, checked identity of Haar
integrals.\chk{where the identity lives: after every private path is integrated, the straight and the L chain leave the same formal word in the two shared links, coefficient for coefficient, in every (sequence, channel) term; after the middle face's private links alone they do not}
\end{theorem}

\begin{remark}[what the lemma does not give]
The lemma is stated and verified for $a,b\le2$. Sixth order reaches $a,b\le3$,
and the $(2,2)$ case and the larger cut families are sampled rather than
exhausted; an inductive proof over $a,b$ would move the mechanism from
``checked at fourth order'' to ``proved at every order''. The corresponding
statement for the single-contact dressing and the fan is verified at the level
of histories but is not covered by this lemma. And the exhaustion shows the
identity is true on the family and false outside it; it does not yet say in one
line \emph{why}.
\end{remark}

\section{Closed forms at every rank}\label{sec:qn}

\subsection{The engine, and running it over $\Q(N)$}

The cumulants are computed by a Wilson-loop calculus written from scratch and
sharing no primitive with the pinned per-rank engine: $H_0$ by the Fierz
identity as a rewiring of index ports (a like pair swaps wires, an unlike pair
is cut and crossed, which is unitarity); Haar integrals by the $U(N)$
Weingarten function realised as the Gram pseudoinverse on $S_n$, with $\det
U=1$ inserted as $\varepsilon\bar\varepsilon$ on virtual slots for the
determinant families; and the resolvent by per-link irrep projectors, so every
component is an exact $H_0$ eigenvector with a known energy.

Nothing in the engine depends on the rank except through rational arithmetic in
$\CF=(N^2-1)/(2N)$, $1/(2N)$ and $N^{\#\mathrm{closed}}$, and through two rank
predicates ($\text{charge}\equiv0\bmod N$ and $|\mathrm{diff}|\le N$).
Promoting those predicates to parameters lets the whole calculation run over
the rational function field $\Q(N)$: scalars become elements of $\Q(N)$, the
Weingarten function becomes the symbolic inverse of the Gram matrix
$N^{\#\mathrm{cycles}}$ on $S_n$, and the single-link spectra become the
$SU(N)$ quadratic Casimirs of the irreps a link's flux content can carry.

\begin{theorem}[the closed forms, and how far the derivation reaches]\label{thm:closed}
Run over $\Q(N)$, each cumulant returns \emph{one} rational function of $N$:
the second-order hops, the two-hop weight, the single-contact and fan dressings
of both pairs, the corner dressing and both cube completions, in both
$C$-parity sectors --- every fourth-order cumulant except the pair cluster,
which is not computed over $\Q(N)$ and does not need to be
(Lemma~\ref{lem:pair}). Each specialises to the pinned per-rank value at every
rank of both records, $N=3,\dots,70$. Every eigencomponent of every resolvent
($21{,}556$ of them, with $103$ distinct intermediate energies) is verified to
be an exact eigenvector of every single-link $H_0$ over $\Q(N)$.

The derivation argument reaches $N\ge5$. At $N=3$ and $N=4$ the forms are
regular and \emph{agree} with the per-rank engine as a checked fact, not as a
consequence: at $N=4$ a resolvent denominator vanishes
(Proposition~\ref{prop:faithful}) and the closed forms additionally rest on
Hypothesis~\ref{hyp:resolvent}'s unaudited determinant trick. The correct
summary is \emph{derived for $N\ge5$, checked at $N=3$ and
$N=4$}.\chk{the specialisation argument's certificate: every eigencomponent an exact H0 eigenvector over Q(N), every Haar family balanced with n <= 3, every flux audited, and the one resolvent denominator vanishing at an integer rank >= 4 is at N = 4}
\end{theorem}

\begin{remark}[one engine, not three]\label{rem:oneengine}
The independence of the $N=3$ corner value (Theorem~\ref{thm:c2}) does not
extend upward. The symbolic-$\Q(N)$ run is the from-scratch calculus with its
scalars replaced by elements of $\Q(N)$: it is the \emph{same} code base. So
every all-rank statement in this section flows through one implementation, and
at $N\ge4$ exactly one exists. Its conventions were moreover matched to the
recorded runs so that its numbers compare directly with theirs, which is what
makes the $N=3$ agreement useful and also what makes it convention-matched
rather than blind. The cheapest thing that would change this is an independent
recomputation of the corner dressing at $N=5$ --- by direct Clebsch--Gordan
contraction, or an orientation-resolved Weingarten evaluator --- against the
pinned rank-sweep row.
\end{remark}

The explicit forms are collected in Appendix~\ref{app:closed}. Three of them
carry the assembly:
\begin{align}
  u(N)&=\frac{N^3(N^2-4)^2\,p_{14}(N)}
             {2(N^2-1)^3(4N^2-9)^3(9N^2-16)(2N^2-3)(2N^2-1)^3(3N^2-2)(4N^2-N-4)(4N^2+N-4)},
  \label{eq:u}\\[2pt]
  d_{\mathrm{single}}(N)&=\frac{2N^3(N^2-4)(10N^2-13)}{(N^2-1)^3(4N^2-9)^2(2N^2-1)^2},
  \label{eq:dsingle}\\[2pt]
  d_{\mathrm{corner}}(N)&=\frac{16N^3(N^2-4)\,p_{18}(N)}
      {(N^2-1)^3(4N^2-9)^3(4N^2-1)(9N^2-25)(2N^2-1)^3(3N^2-1)(4N^2-5)(4N^2-2N-5)(4N^2+2N-5)},
  \label{eq:dcorner}
\end{align}
with $p_{14}$ and $p_{18}$ the even polynomials given in
Appendix~\ref{app:closed}. The double zero $(N^2-4)^2$ in the C-odd two-hop
weight, and the simple zero $(N^2-4)$ in the dressings, are the same $SU(2)$
degeneration that makes $t_N$ vanish at $N=2$.

Every denominator factor is an energy denominator $E_0-E_j$ of an identified
intermediate state, and this is checked rather than asserted: the corner adds
$3N^2-1$ from $(2,\adj\adj)$, $2N\mp1$ from $(2,\adj\adj\adj)$,
$4N^2\mp2N-5$ from $(2,\adj\Lambda\Lambda)$ and $(2,\adj\Sym\Sym)$, $4N^2-5$
from $(2,\adj\Lambda\Sym)$, and $3N\mp5$ from $(2,\Lambda\Lambda)$ and
$(2,\Sym\Sym)$.\chk{the corner dressing by channel: 87 C-odd channels, at most three doubled links in any direct-term state, every resolvent factor the corner adds to R20 the energy of an identified intermediate state, and the Weingarten poles at N = -+3 of single channels cancelling in the sum}

\begin{proposition}[the cube completions]\label{prop:cubes}
Every cube link lies in exactly two faces, so once both are present the link's
adjoint component can never return to the target and only the singlet --- the
$n=1$ Haar moment --- survives. Extending that primitive law to multi-loop
intermediates gives, at every rank,
\begin{equation}\label{eq:cubes}
  K_{\mathrm{opposite}}(N)=\frac{-160}{N(N^2-1)^3},\qquad
  K_{\mathrm{adjacent}}(N)=\frac{-106}{N(N^2-1)^3},
\end{equation}
ratio $53/80$ at every $N$. The adjacent value is the primitive $-88$ (the
off-axis ledger's $S_4=-11$) plus ten multi-loop orderings at $-18$. At $N=3$,
$K_{\mathrm{adjacent}}=-53/768$.\chk{the adjacent-face cube completion is -106/(N(N\^{}2-1)\^{}3): the primitive -88 plus ten multi-loop histories at -18}
\end{proposition}

\subsection{The assembly collapses to three cumulants}

\begin{proposition}[coplanar equals minus perpendicular]\label{prop:coplanar}
As rational functions of $N$, the coplanar single-contact and fan dressings are
the exact negatives of the perpendicular ones in the C-odd sector, and equal to
them in the C-even sector. With the pair cluster cancelling
(Lemma~\ref{lem:pair}), the fans drop out of $\rho+\pi$ entirely, which is why
$R_{20}$ carries none of the fan's factors
$(4N\mp3)(4N\mp7)(4N^2-7)$.\chk{the coplanar dressings are the exact negatives of the perpendicular ones in Q(N), single contact and fan alike, so the eleven-cumulant assembly is three cumulants and the cube: beta\_N = -16u + 32d - 16 corner + 848/(N(N\^{}2-1)\^{}3)}
\end{proposition}

\begin{lemma}[the pair cluster cancels, by parity]\label{lem:pair}
The coplanar and perpendicular two-face clusters are \emph{one} abstract graph
up to conjugating a single face. Every link bijection carrying coplanar words
onto perpendicular words conjugates an odd number of faces --- the parity set
found is $\{1\}$ --- so in the C-odd sector the two are exact negatives at
every rank and the stacked pair is identically zero. In particular the pair
cluster's double pole at $N=3$, the $(3,0)$ determinant family, is invisible to
the shape coefficient.\chk{the coplanar and perpendicular two-face clusters are one abstract graph up to conjugating Q, so their C-odd pair elements are exact negatives at every rank and the pair cluster drops out of C\_shp}
\end{lemma}

\begin{remark}[the pair cluster is the one that is never computed]
This lemma is load-bearing and its status should be stated plainly. The pair
cluster is the single fourth-order cumulant with \emph{no} closed form: it is
absent from the $\Q(N)$ record entirely, it is null in the $N=4$ assembly
certificate, and it is infeasible at $N=4$. It costs nothing only because it
cancels --- and the cancellation is an abstract-graph parity argument, not a
computation. So ``$\beta_N$ from three cumulants'' is exact, but ``the
eleven-cumulant assembly was computed'' would not be: ten of eleven were, and
the eleventh was argued away.
\end{remark}

\begin{theorem}[$\beta_N$ from three cumulants]\label{thm:beta}
Let $\alpha$ be the eleven-cumulant assembly's cube amplitude. Then $\alpha$
cancels identically:
\[
  8\Bigl(\frac{\alpha}{4}\Bigr)
  +16\Bigl(-\frac{\alpha}{8}-u-\frac{(\rho+\tilde\pi)}{2}\Bigr)
  \;=\;-16u+32\,d_{\mathrm{single}}-16\,d_{\mathrm{corner}}-8K ,
\]
and with $K=K_{\mathrm{adjacent}}=-106/[N(N^2-1)^3]$,
\begin{equation}\label{eq:betaid}
  -16\,u(N)+32\,d_{\mathrm{single}}(N)-16\,d_{\mathrm{corner}}(N)
  +\frac{848}{N(N^2-1)^3}
  \;=\;\frac{P_{17}(N^2)}{N\,R_{20}(N^2)}
\end{equation}
wherever $N\,R_{20}(N^2)\ne0$, where $P_{17}$ and $R_{20}$ are the corpus's
degree-$17$ numerator and degree-$20$ factored denominator in $z=N^2$
(Appendix~\ref{app:p17}). Every cumulant denominator divides $NR_{20}(N^2)$;
over that common denominator the weighted numerators sum to $P_{17}(N^2)$
exactly --- a polynomial identity of degree $34$ in $N$. The identity, the four
cofactor lemmas and the assembly are proof-checked in Lean~4
(\texttt{betaN\_from\_three\_cumulants}, \texttt{betaN\_assembled\_numerator},
\texttt{assembly\_three\_cumulants}), with no \texttt{sorry} and only
\texttt{propext}, \texttt{Classical.choice} and \texttt{Quot.sound}.\chk{the final polynomial identity: -16 u + 32 d - 16 corner - 8 K\_adj = P17(N\^{}2)/(N R20(N\^{}2)) as rational functions, for the closed forms of the two-hop weight, the single-contact dressing, the corner dressing and the adjacent-face cube completion}
\end{theorem}

\begin{remark}[what the degree bound was, and why it is gone]
The same equality had been established by assembling the eleven cumulants at
$N=4,\dots,70$ and observing agreement at every rank. That gives the identity
for all $N$ only \emph{conditional} on $\beta$ being rational of degree at most
$(17,20)$ in $N^2$ --- $67$ exact agreements plus a plausibility argument. One
could have argued the bound from the engine's structure (Weingarten
denominators, energy denominators, a vertex count) and then invoked uniqueness
of rational interpolation. That argument was not made. Theorem~\ref{thm:closed}
makes it unnecessary: the degree of $\beta_{\mathrm{assembled}}$ is now the
\emph{computed} degree, because the sum of derived rational functions is a
derived rational function.
\end{remark}

\begin{remark}[what the identity certifies, and what it does not]\label{rem:twopipelines}
Theorem~\ref{thm:beta} is an agreement between two pipelines: the corpus's
generator, which produced $P_{17}/R_{20}$ and is \emph{output-certified only}
--- its generator chain was never regenerated end to end here, and the
certificate for its own internal reduction is absent from this repository ---
and the engine of Theorem~\ref{thm:closed}, whose premises are
Hypotheses~\ref{hyp:pvp}--\ref{hyp:resolvent}. Agreement of this strength
between two independently written pipelines is evidence for both. It is not
proof of either, and the cold one-shot regeneration of the all-rank
fourth-order bundle is an open item. What Lean certifies is narrower and
unconditional: that the four displayed closed forms, with weights
$(-16,32,-16,-8)$, sum to the displayed $P_{17}(N^2)$ over the common
denominator. That those forms \emph{are} the engine's cumulants is the exact
check beside it; that the engine's cumulants are the physics is
Hypotheses~\ref{hyp:pvp}--\ref{hyp:resolvent}.
\end{remark}

\begin{proposition}[faithfulness of the specialisation]\label{prop:faithful}
The symbolic and per-rank engines run the same word arithmetic and can differ
only through the rank predicates, the Weingarten function and the spectral
decomposition. The record certifies: no word ever carries every nonzero link
flux at magnitude $\ge5$ (largest flux $6$, and every such word also carries one
$\le4$), so the two charge filters agree for $N\ge5$; every Haar family is
balanced with $n\le3$, so the per-rank pseudoinverse is the inverse, which is
the specialisation of the symbolic one. The resolvent denominators $E_0-E$ have
exactly one integer root at $N\ge4$: at $N=4$, from the state $P^3$ with every
link of the plaquette in $\Lambda^3F$, which is $\bar F$ at $N=4$, so the state
is degenerate with the plaquette. Hence for $N\ge5$ the closed forms are the
engine's values by argument; at $N=3$ and $N=4$ they are regular and agree with
the engine as a checked fact.\chk{every cumulant of the beta\_N assembly is one rational function of N computed over Q(N): the closed forms are derived, not reconstructed, and specialise to the per-rank records at N = 3..70}
\end{proposition}

\begin{corollary}[$\beta_N$ is a sum of resolvent products]\label{cor:resolvent}
$\beta_N$ is a sum of $254$ channel terms with the assembly's integer weights,
every denominator a product of the resolvent factors above (together with $N$,
$N\pm1$, $N\pm3$) of degree at most $7$, every numerator of degree at most $4$,
summing to $P_{17}(N^2)/[N R_{20}(N^2)]$ identically. Single channels carry the
$n=3$ Weingarten poles $N=\pm3$; the sum does not.\chk{beta\_N is a sum of resolvent products: -16 u + 32 d - 16 corner + 848/(N(N\^{}2-1)\^{}3) over 254 channel terms, each a rational function of N whose denominator is a product of resolvent factors, equals P17(N\^{}2)/(N R20(N\^{}2))}
\end{corollary}

\section{The off-axis coefficient}\label{sec:c2}

\subsection{What the dispute was, and what it is not}

Two fourth-order kernels were recorded, with off-axis coefficients
$C_{\mathrm{old}}=-0.0480863\ldots$ and $C_{\mathrm{new}}=-0.0202133\ldots$.
The first thing to establish is that no re-anchoring can reconcile them.

\begin{theorem}[the crosswalk cannot decide it]\label{thm:obstruction}
The two coordinate systems are related by
\[
  c_4^{\mathrm{new}}(k)=c_4^{\mathrm{old}}(k)+\Delta_\Gamma
    +\Delta_C\,\Phi_C(k),\qquad \Phi_C(k)=\frac{4e_2(k)}{q(k)} ,
\]
and $\Phi_C(0)=0$. So the $\Gamma$-point scalar pins $\Delta_\Gamma$ and places
\emph{no} constraint on $\Delta_C$. $\Phi_C$ also vanishes on every axial cut,
which is exactly why axial data agree while $M$ and $R$ split by $8\Delta_C$
and $16\Delta_C$. No $\Gamma$ or axial datum can identify the off-axis
coefficient; an explicit second witness exhibits the non-identifiability by
construction.\chk{FINDING: the retained Gamma/axis data cannot identify C\_shp}
\end{theorem}

\subsection{Solved, not fitted}

The recorded comparison fitted the ansatz at four points of the zone with
condition number $45.8$, which is why every non-axial number in the dispute was
a float.

\begin{theorem}[exact solution over the whole zone]\label{thm:solved}
Clearing the $1/q$ turns \eqref{eq:ansatz} into a linear system over Laurent
polynomials in $z_j=e^{ik_j}$, which has an exact solution over the whole zone
or none at all. It has one:
\[
  A=\tfrac{5}{48},\qquad B=0,\qquad D=0,
\]
with residual Laurent polynomial identically zero, and $C$ equal to the
registry's recorded value. The four-shape form is therefore not an artefact of
where the transcript sampled.\chk{the shape coefficients are solved, not fitted: A = 5/48, B = D = 0 exactly}
\end{theorem}

\begin{corollary}[the dispute is one scalar]\label{cor:onescalar}
In the orbit basis of Proposition~\ref{prop:orbits}, $A=5/48$ forces the normal
amplitude $\nu=-(5/48+4u)$, and the on-site amplitude $\sigma$ enters $c_0$
alone. The whole of the dispute is
\begin{equation}\label{eq:cshpform}
  C_{\mathrm{shp}}=-\tfrac{5}{96}-u-\tfrac{1}{2}(\rho+\tilde\pi) ,
\end{equation}
three signed numbers. The recorded kernels have the same six orbits and the
same normalised table row for row; the second one's skeleton unit is
$4.1327437\times$ and its $\rho$ and $\tilde\pi$ are \emph{opposite in
sign}.\chk{FINDING: the cold kernel is the same six orbits with rho and pi sign-flipped}
\end{corollary}

Since $\tilde\nu$ agrees between the kernels, no normalisation relates them: the
$4.13\times$ scale on $u$ is a genuine fourth disagreement, and it is
band-invisible on the carrier because $u$ reaches the band only through $c_0$.

\subsection{The decision}

\begin{theorem}[the off-axis coefficient]\label{thm:c2}
At $N=3$, three implementations of the corner cluster --- the pinned engine's
run, a Wilson-loop calculus written from scratch that shares no primitive with
it, and the original pipeline's own recovered word ledger --- agree to the
digit:
$-2580244782961/398756546697600$ in the C-odd sector and
$-56022878647/4153714028100$ in the C-even, with all four orientation entries
identical. The original pipeline's own face-resolved word ledger --- $4{,}221$
ordered words, recovered from the corpus, whose hash matches the one the kernel
copies quote, and whose rooted assembly reproduces the pinned $189$ records
exactly --- agrees with the cluster assembly on \emph{every} cluster of the
rotation record (pair, fourteen chains, two fans, two corners: identical
rationals in both sectors) except one: the adjacent-face cube completion, where
it holds $-31/1536$ against $-53/768$. It carries $8$ of the $24$ insertion
orderings, each with exactly the weight the multi-loop law assigns it
($2,1,4,2,8,8,2,4$ in units of $1/1536$); the sixteen it lacks sum to
$-25/512$. The opposite-face completion in the same ledger has all $24$.
Consequently, in the kernel's basis,
\[
  \rho_{\mathrm{assembled}}=-\frac{588708011765248393}{14501180577204921600},
  \qquad
  \rho_{\mathrm{historical}}-\rho_{\mathrm{assembled}}=\frac{25}{512},
\]
\[
  C_{\mathrm{shp}}=C_{\mathrm{historical}}+\frac{25}{1024}
  =-\frac{13035490122347}{550663802582400}.
\]
Put into the $24$ records and solved over the whole zone: $A=5/48$, $B=D=0$, no
residual. The arithmetic chain --- the shortfall $53/768-31/1536=25/512$, its
propagation $C\mapsto C+25/1024$ under $\rho\mapsto\rho-25/512$, the resulting
value, and the induced shift of $\beta$ --- is proof-checked in
Lean.\chk{FINDING: C\_shp from the assembled amplitudes in the kernel's basis is C\_historical + 25/1024, the registered continuation-shifted value}
\end{theorem}

\begin{remark}[adjudicated, and the falsifier that is unperformed]\label{rem:adjudicated}
Theorem~\ref{thm:c2} is an adjudication between two recorded kernels, using a
third engine's cube-ordering enumeration and the original pipeline's own
ledger. It is not a measurement of the coefficient from first principles. The
repository's own declared falsifier is named and has not been performed: no
direct balanced contraction of the fourth-order kernel at $N=3$ is held here.
Until that exists, the honest statement is that the contradiction is
\emph{adjudicated} --- every component but one is the original pipeline's own
number, and the one is explained ordering by ordering --- and that both
recorded sides remain in the register with their verdicts rather than being
deleted.
\end{remark}

Two corrections of the record come with it. First, the assembled rotation
element was previously read in the conjugate basis: all $24$ cross-plane
records equal $\rho$ times the $x$-then-$z$ incidence, and every C-odd component
of the historical ledger is the negative of the assembly's $(2,0)$-traversal
value while every C-even one is identical --- the signature of conjugating the
$(0,2)$ face and nothing else. Second, the corpus's recorded rule
``determinant sectors shift $C$ by $-25/1024$ at $N=3$'' is this shortfall, on
a cluster that has no determinant sector; the predicted
$\Delta(\rho+\tilde\pi)=-25/512$ is realised as the ledger's shortfall and not
as the $\varepsilon$-sector, which was measured at $-55/6936$.

\begin{remark}[the third recorded value, and why it moved]
Between the two kernels and \eqref{eq:cshp} there was, for two days, a third
value. Substituting the cluster-assembled amplitudes into the orbit-shape
formula \eqref{eq:cshpform} gave
$C^{\mathrm{cl}}_{\mathrm{shp}}=-54822624038066723/853010622188524800
=-0.0642696\ldots$, recorded beside the historical $-0.0480864$ and the second
kernel's $-0.0202133$ and promoted no further than they were, together with the
statement that the calculation which would close the question was the corner
cluster from a third implementation. That calculation was then done. The
assembly was right; its \emph{basis} was not --- $\rho$ had been read in the
conjugate basis --- and correcting the basis moves the third value to
\eqref{eq:cshp}. The intermediate value is retained in the record with its
verdict, and the check that produced it still verifies its own arithmetic and
says so.
\end{remark}

\begin{remark}[a forbidden route that reached the right number]
The value \eqref{eq:cshp} coincides with what an analytic continuation of the
all-rank formula to $N=3$ produces. That continuation was retracted as a
\emph{derivation}, because the formula has a third-order pole there, and the
retraction stands. What Theorem~\ref{thm:c2} establishes is that the number the
forbidden route produced is the number the cluster assembly reaches
independently. The distinction is the point: a value is not promoted because a
route reached it, and both recorded sides remain in the registry with their
verdicts.
\end{remark}

\section{The $SU(3)$ shape row}\label{sec:su3}

\begin{corollary}[what moves, and by how much]\label{cor:table}
By Theorem~\ref{thm:obstruction}'s crosswalk, a change $\Delta_C$ moves nothing
on the axis, $8\Delta_C$ at $M$ and $16\Delta_C$ at $R$, hence $16\Delta_C$ on
$\beta=8A+16C$ and on the bandwidth $W_4=\alpha+\beta$. With
$\Delta_C=25/1024$:
\[
\begin{array}{lll}
  \alpha=4A=\tfrac{5}{12} & \text{unchanged} & \lambda_X \text{ unchanged}\\[2pt]
  \beta_3: & \dfrac{17607806155349}{275331901291200}\;\longrightarrow\;
             \dfrac{15644916262153}{34416487661400} & (+\,25/64)\\[8pt]
  W_4: & \dfrac{132329431693349}{275331901291200}\;\longrightarrow\;
         \dfrac{29985119454403}{34416487661400} & (+\,25/64)\\[8pt]
  \lambda_M=\alpha+\beta/2: & & (+\,25/128)\\[2pt]
  \lambda_R=W_4: & & (+\,25/64)
\end{array}
\]
The near-$\Gamma$ exclusion radius, proportional to $\sqrt{W_4}$, grows by the
corresponding factor.\chk{the fourth-order shape table with the resolved C\_shp: beta\_3, W\_4, lambda\_M and lambda\_R move by 25/64, 25/64, 25/128 and 25/64; alpha and lambda\_X do not}
\end{corollary}

\begin{corollary}[the $SU(3)$ exception was an enumeration shortfall]\label{cor:noexception}
The assembly of Theorem~\ref{thm:beta} at $N=3$ gives
$8A+16C_{\mathrm{shp}}=15644916262153/34416487661400$ with $A=5/48$ and
$C_{\mathrm{shp}}$ as in \eqref{eq:cshp}, and this is exactly the number the
corpus's compact formula produces at $N=3$. The recorded ``$SU(3)$ exception in
the shape channel'' was therefore the sixteen missing cube orderings, and the
recorded rule attributing it to determinant sectors is retracted as a
mechanism: the shifting cluster has no determinant sector at all --- every link
carries $U\otimes U\dg$ --- and the $\varepsilon$ sector proper was measured at
$-55/6936$, not the predicted $-25/512$.\chk{the all-rank beta\_N formula at N = 3 is 8A + 16 C\_shp\_assembled exactly: the SU(3) exception in the shape channel was the cube shortfall}
\end{corollary}

\begin{remark}[the formula is still not a derivation at $N=3$]
Corollary~\ref{cor:noexception} is a statement about two numbers agreeing, not
a licence to evaluate the compact formula at $N=3$. That substitution stays
forbidden: the continuation route has a third-order pole there. The two
rationals are deliberately kept in the registry under separate names, and the
derivation at $N=3$ is the assembly, never the formula.
\end{remark}

For completeness, the six orbit amplitudes at $N=3$ in the kernel's basis are
in Appendix~\ref{app:ledger}, together with $\tilde\nu=-5/48$, the axial
all-rank cube coefficient $\alpha_N=640/[N(N^2-1)^3]$, and
$u(3)=360421351/40327601932800$ with $u_2=2u$ exactly.

\section{One angular invariant: the anisotropy variance}\label{sec:aniso}

Section~\ref{sec:collapse} showed that leaving the carrier costs one $R$.
This section computes what that costs, and finds that the whole of it is a
single scalar function of the direction of $k$.

Write $x_i=a_i/q$, so $x_1+x_2+x_3=1$ on the simplex, and let $p=\psi/\sqrt q$
be the normalised carrier ($q>0$; $\Gamma$ is excluded throughout this
section, where $q=0$ and $\psi=0$). Let $H=H_4-\tilde\sigma I$ be the centred
kernel.

\begin{theorem}[the mixing is one invariant]\label{thm:aniso}
With
\[
  V(x)=\sum_i x_i^3-\Bigl(\sum_i x_i^2\Bigr)^{2}
      =\sum_{i<j}x_ix_j(x_i-x_j)^2 ,
\]
one has $QHp=-2C_{\mathrm{shp}}\,QDp$ and
\begin{equation}\label{eq:aniso}
  \bigl\|QHp\bigr\|^2 \;=\; 4C_{\mathrm{shp}}^2\,q^2\,V(x) .
\end{equation}
This is verified as an exact cleared Laurent identity: all $85$ nonzero
Laurent coefficients of $q\,\langle\psi,H^2\psi\rangle-\langle\psi,H\psi\rangle^2$
agree with $4C_{\mathrm{shp}}^2\sum_{i<j}a_ia_j(a_i-a_j)^2$, with no
tolerance. The underlying shape identity $V=e_2+3e_3-4e_2^2$ on
$x_1+x_2+x_3=1$ is proof-checked in Lean.\chk{anisotropy mixing residual matches the full assembled Laurent kernel}
\end{theorem}

\begin{proposition}[nodal classification]\label{prop:nodal}
$V\ge0$ on the nonnegative octant, with $V=0$ exactly when every pair of
coordinates has a zero or is equal --- that is, when all positive $a_i$ are
equal: the axes, the equal face diagonals, and the equal body diagonals. On
the simplex $V\le1/4$. A generic holdout: $V(1/6,1/3,1/2)=5/324>0$. Both the
nonnegativity and the $1/4$ bound are proof-checked in Lean.\chk{anisotropy variance is exactly the simplex pair sum and cubic shape invariant}
\end{proposition}

So the carrier is an exact eigenvector of the fourth-order kernel precisely on
the isotropic directions, and the deviation from flatness is measured by a
variance of the lattice $\sin^2(k_i/2)$ about their mean. Equality here is
equality of the $a_i$, not of the $k_i$.

\begin{proposition}[the induced sixth-order coefficient]\label{prop:sixth}
Nondegenerate perturbation theory on the displayed truncated matrix
$M/u^2=t_3\,q\,Q+u^2H$ gives
\[
  E_0=u^4\,p\dg H p-\frac{4C_{\mathrm{shp}}^2}{t_3}\,u^6\,q\,V+O(u^8),
\]
so the fourth-order virtual mixing induces
\[
  \delta E_{6,\mathrm{mix}}
  =-\frac{4C_{\mathrm{shp}}^2}{t_3}\,u^6
    \Bigl(\frac{e_2}{q}+\frac{3e_3}{q^2}-\frac{4e_2^2}{q^3}\Bigr),
\]
a definite negative coefficient with shape ratio exactly $1:3:-4$. With
$t_3=5/612$ and $C_{\mathrm{shp}}$ as in \eqref{eq:cshp},
\[
  -\frac{4C_{\mathrm{shp}}^2}{t_3}
  =-\frac{169924002729806205028788409}{619343593697933825385600000}
  =-0.2743614440495553\ldots
\]
The rational identity is proof-checked in Lean; the operator origin is
not.\chk{anisotropy induced sixth-order coefficient is minus four C squared over t3}
\end{proposition}

\begin{corollary}[the branch is decidable at sixth order]\label{cor:branch}
The coefficient depends on $C_{\mathrm{shp}}$ only through $C_{\mathrm{shp}}^2$.
Substituting the historical branch instead of \eqref{eq:cshp} changes it by a
factor of about $4.1$. So the decision of Section~\ref{sec:c2} is not internal
bookkeeping: it fixes a sixth-order number, and a sixth-order census is in
principle a test of it.
\end{corollary}

\begin{remark}[what this is not]
$\delta E_{6,\mathrm{mix}}$ is the induced-mixing contribution of the displayed
truncated matrix at fixed nonzero momentum. It is not the complete sixth-order
coefficient: a genuine sixth-order matrix contributes its own carrier
expectation, and fifth-order terms and other sectors belong as well. Whether
the $1:3:-4$ shape survives is an explicit open calculation. The identity
\eqref{eq:aniso} also yields overlap and energy bounds ---
$1-|p\dg v_0|^2\le 4C_{\mathrm{shp}}^2u^4V/d^2$ and
$0\le\mu-E_0\le4C_{\mathrm{shp}}^2u^6qV/d$ with $d=t_3-2\cdot\tfrac{5}{48}u^2>0$
for $0<u^2<2/51$ --- and the overlap bound carries no $q$ in the denominator,
hence no near-$\Gamma$ volume loss. These control the carrier's overlap with
the lowest band of a finite-order truncated matrix. They are not the physical
source overlap.
\end{remark}

\section{External comparisons}\label{sec:external}

Independent published work matters here because it was produced without
knowledge of this program, not because it is published. Three comparisons are
exact, one is a reconciled apparent disagreement, and one names the same
obstruction from the other side.

\paragraph{Exact, rational for rational, at five orders.} Kogut, Pearson and
Shigemitsu's 1980 table of $SU(3)$ Hamiltonian string-tension coefficients,
bridged by their own printed coupling convention ($\sigma_n=2^{n-1}t_n$ with
$x=2u$), equals this program's independently computed series \emph{exactly} at
$n=0,2,3,4,5$ --- including an $18$-digit numerator at $n=4$ and a $24$-digit
CRT-reconstructed numerator at $n=5$:
\[
  \sigma_4=-\frac{737327120374220449}{7250590288602460800},\qquad
  \sigma_5=-\frac{137767222189182735950309}{2009803206414863779920000},
\]
with no tolerance anywhere. The $n=5$ row agrees with the authors' own 1980
correction of their 1979 error, not with the original. Their $t_1=0$ matches
the absent $\sigma_1$, and the uniform negativity of their printed table
settles this program's sign convention independently. This is the strongest
external anchor the program has, and it is exact. It is not, however, a cold
recomputation: $\sigma_5$ here is a seven-prime CRT reconstruction whose native
rerun is an open item, and $\sigma_4$ shares a denominator with the
fourth-order kernel lineage and inherits whatever that
settles.\chk{the KPS 1980 string-tension table equals the certified sigma series EXACTLY}

\paragraph{Order by order against a 1989 glueball table.} Hamer's $SU(3)$
strong-coupling glueball series agrees with this program's vacuum-subtracted
$\Gamma$-point coefficients at every shared order through $x^3$, in both the
$1^{+-}$ and $0^{++}$ channels, with the sign structure matching ($a_1=-1$ in
the scalar channel where the carrier's is $+1$). These are float comparisons
against printed twelve-digit decimals, at stated tolerances --- T2, not exact,
and they should not be quoted alongside the previous paragraph as though they
were. The bridge $m_n=2^{n-1}a_n$ is proved as algebra from Hamer's own printed
equations, not fitted, so no $4^r$ ambiguity enters.\chk{the m\_n = 2\^{}(n-1) a\_n bridge is the x = 2u conversion}

\paragraph{A published truncation that disagrees, and why.} The link cutoff
$p+q\le1$ used for published one-cube exact diagonalisations gives a
charge-odd hopping \emph{opposite in sign} to the channel-complete value and
$10.2\times$ larger in magnitude. This is not a discrepancy in the physics: by
arithmetic on the weight ledger alone, a cutoff keeping only the singlet and
$\Lambda^2F$ routes projects the hopping to $-1/12$ and omits exactly $14/153$,
the two summing to $t_3=5/612$. Rebuilding a $\mathcal B=4$ open cube from one
such group's own public plaquette table reproduces $-1/12$ as a matrix element
together with the reversed shell ordering; the $\mathcal B=6$ cube flips the
sign back, and $t_3$ is then the off-diagonal element of the gap operator with
the shell ordering exactly inverted. The disagreement is a property of the
truncation, and it is reconciled channel by channel.\chk{FINDING: the B = 6 cube flips the sign back to +5/612, closing the decisive test}

\paragraph{The same obstruction, from the Euclidean side.} Schierholz's
$SU(3)$ glueball paper states independently the exact obstruction this program
records at $G18$: the bare local operator's overlap with the physical state
degrades with the lattice spacing. Section~\ref{sec:spectral} explains that
structurally rather than as an empirical constant --- the exactly flat local
algebra cannot see the state at all, so the coupling is entirely
dressing-borne.\chk{the overlap obstruction was published in 1988, and it scales}

\paragraph{The recorded external disagreement.} One published claim contradicts
this program's values, at $x^3$ and $x^4$, and it is recorded as such rather
than resolved: it is a second-hand report about a 1976 three-author paper, and
the primary source has not been obtained. It is registered with that standing.
Provenance was disentangled in the process: the series in question is
Kogut--Sinclair--Susskind 1976, not the two-author Kogut--Susskind 1975
Hamiltonian paper.\chk{the series Hamer supersedes is Kogut-SINCLAIR-Susskind 1976, not KS 1975}

\section{What the fourth order buys at $\Gamma$ and near it}\label{sec:gamma}

\begin{proposition}[the $\Gamma$ triplet is a symmetry fact]\label{prop:gamma}
At $\Gamma$ the three plaquette orientations carry one irreducible $T_1$
representation of the cubic group on the three-dimensional orientation fibre.
By Schur's lemma, any operator on \emph{that fibre} which commutes with the
cubic action is a scalar there, at every order in $u$; and inversion kills the
linear term. So the triple degeneracy at $\Gamma$ is exact and is a spin-one
multiplet at rest, rather than a coincidence of the displayed
coefficients.\chk{the three plaquette orientations at k = 0 carry an irreducible representation of the cubic group (the T1 triplet), so every cubic-covariant effective Hamiltonian is a scalar at Gamma at every order}
\end{proposition}

The restriction matters: this is a representation-theoretic statement about a
$3\times3$ block, silent about states outside the fibre and about whether the
effective Hamiltonian exists as an operator at all orders.

\begin{proposition}[fourth-order isolation away from $\Gamma$]\label{prop:iso}
On the fourth-order truncation $M(k)=t(u)\Lambda(k)+u^4\bigl(H_4(k)-\tilde\sigma
I\bigr)$ one has the tight bound $\|H_4(k)-\tilde\sigma I\|\le\tfrac{5}{48}\,q_a(k)$,
attained on an axis, and hence by Weyl's inequality the lowest band of $M$ is
isolated at \emph{every} $k\ne0$ for $u^2<2/51$ --- with no exclusion radius,
which supersedes the earlier sufficient criterion $|k|\ge Ku$.\chk{the fourth-order kernel is bounded by C\_iso q\_a(k) away from its Gamma scalar, C\_iso = 5/48 for the assembled C\_shp, so the flat band is isolated at every k != 0 for u\^{}2 < 2/51 -- no exclusion radius}
\end{proposition}

Three qualifications, all of them load-bearing. This is a statement about the
fourth-order \emph{truncation}: every order $\ge5$ is dropped, and bounding
their sum uniformly is exactly the open problem. The threshold $u^2<2/51$ uses
$t(u)\ge t_3u^2$, which needs the positivity of an upstream third-order number
whose raw chain is not replayed here. And the isotropy constant is the same
$5/48$ that Corollary~\ref{cor:carrier} produced, which is why this proposition
belongs to the fourth order rather than to the band.

The sharp zone minimum, used wherever a momentum ball is excluded, is
$\min_{|k|\ge r}q=4\sin^2(r/2)$; it is uniform in $L$ and should not be
confused with the minimal-grid-momentum statement $4\sin^2(\pi/L)$, which
concerns a different set.\chk{the exact zone minimum of q on |k| >= r is 4 sin\^{}2(r/2), and Jordan overstates it by pi/2}

\section{Spectral standing at fixed lattice spacing}\label{sec:spectral}

The results above concern an effective operator on a projected shell. This
section records, without reproving, how far that operator has been connected to
a spectral statement, because the boundary is where overclaims live.

\paragraph{Finite volume, all orders.} A local spectral enclosure of Yarotsky
type gives, for $SU(3)$, an all-orders finite-volume Riesz island: the complete
three-component cluster persists on an existential small-$|u|$ domain, with
constants uniform in $L$. This supplies an interacting Riesz cluster, not a
particle; it does not prove convergence of the displayed coefficient series,
persistence of one internal sheet, or identification with the lightest physical
glueball.

\paragraph{Infinite volume, fixed spacing.} The thermodynamic ground state and
infinite-volume Hamiltonian come from the same quantum-perturbation input
applied to a three-link cell decomposition; a close-projection and GNS-transport
chain carries the coefficient shell into the physical sector, and the
second-order coefficient of the exact orthonormal symbol isolates a rank-one
internal sheet on every momentum ball compactly separated from $\Gamma$. Two
boundaries survive: no rank-one separation is claimed \emph{at} $\Gamma$, where
the cubic $T_1$ triplet meets; and the identification of the resulting
semigroup with an isotropic Wilson-action transfer matrix is a separate
problem.

\paragraph{Near $\Gamma$.} At fixed nonzero momentum the in-shell splitting
$t_Nu^2q(k)$ is $L$-independent along nested volumes, and the shell around the
carrier keeps a volume-uniform electric isolation. The $L^{-2}$ collapse that
older accounts cite is the \emph{minimal-grid-momentum} statement, not a
fixed-momentum one: a particle statement lives at fixed $k$, where no collapse
occurs, and the collapse governs only the approach to $\Gamma$. The sharp zone
minimum is $\min_{|k|\ge r}q=4\sin^2(r/2)$.

\paragraph{The continuum.} Blocked, and two of the obstructions are theorems;
see Section~\ref{sec:notclaim} and Section~\ref{sec:scope}.

\section{Negative results}\label{sec:negative}

A program of this size is better characterised by what it has ruled out than by
what it has assembled. The following are all recorded with explicit witnesses.

\begin{enumerate}\itemsep3pt
\item \textbf{The degree bound (retracted).} The vertex-counting argument for
  the tier collapse is false; witness
  $d\dg\operatorname{diag}(a_i^2)d=q^3-3qe_2+3e_3$. Its prediction that the
  $L^{-4}$ tier appears first at sixth order is withdrawn. See
  Remark~\ref{rem:retraction}.
\item \textbf{The $R$-degree selection rule (refuted).} $\sigma(UR)=-2qe_2$;
  Finding~\ref{find:ur}.
\item \textbf{The elementary-symmetric obstruction hypothesis (refuted at the
  level of the algebra).} $\sigma(RUR)=4e_2^2$; Finding~\ref{find:rur}.
\item \textbf{The $F$-$2$ mobility rule (refuted).} The promotion
  $r_{\mathrm{physical}}=w_{\min}-2$ is falsified by the pentagonal isotropic
  cap hop at $r=4$ with $w_{\min}-2=5$. The scoped bound
  $r\ge w_{\min}-2$ survives, with survival gates.
\item \textbf{The stranded-flux zero backend (refuted).} A claimed Haar zero
  for all $20$ pentagonal histories is refuted by exact Weingarten and network
  contractions.
\item \textbf{Center-only circuits are dynamically dark (refuted).} Falsified
  by the exact fifth-order determinant correction $\delta c_{5,\det}\ne0$.
\item \textbf{The global fixed-window firewall (disproved).} Replaced by the
  rooted conditional calculus.
\item \textbf{Sector monotonicity is dead on a sign, not on an estimate.} A
  whole class of strong-to-weak-coupling transport arguments --- Griffiths-type
  sector-wise monotonicity in $\beta$ --- is closed before the inequality is
  attempted. The leading electric-flux sector gap is $\CF L/(2\sqrt u)$, with
  derivative $-(N^2-1)L/(8Nu^{3/2})$: strictly \emph{decreasing} in $u$, hence
  in $\beta$, at every rank $N\ge2$. And the slope coefficients
  $(n-\tfrac12)\sigma_n$ are negative at every certified order,
  \[
    -\tfrac13,\quad -\tfrac{11}{51},\quad -\tfrac{305}{816},\quad
    -\tfrac{737327120374220449}{2071597225314988800},\quad
    -\tfrac{137767222189182735950309}{446622934758858617760000},
  \]
  so such a monotonicity, if it holds at all, bounds the weak-coupling gap from
  \emph{above} and never from below. The falsifier that would revive the route
  is recorded: a functional of the sector free energies, defined without
  knowledge of the large-$\beta$ behaviour, increasing in $\beta$ at small
  $\beta$, whose positivity at large $\beta$ implies a gap.
\item \textbf{The leak-equals-hop identification (refuted; the equality
  survives).} The equality is derived for a reason that does not involve the
  vacuum route, so the identification falls while the equality stands and both
  rationals keep their registered values. The C-odd control decides which
  family the surviving cross-pair carries: $-3/68$ at $N=3$, not $-1/36$ ---
  which is why the C-odd sector separates at third order and the C-even one
  does not.
\item \textbf{The $(\rho,\pi)$ sign flip is not a convention (refuted).} Under
  the orientation character every one of the six orbits is separately Hermitian
  and invariant under all $48$ elements of the cubic group, so no orientation
  convention relates the two kernels.
\item \textbf{A single weight on $\Lup$ is not the mechanism.} $-160$ against
  $-106$; Proposition~\ref{prop:cubes}.
\item \textbf{One recorded kernel is wrong on $u$, $\pi$ and $\rho$; the other
  on the adjacent-face cube completion.} Theorem~\ref{thm:c2}.
\item \textbf{The stochastic blocking variant (refuted).} Equivariance plus
  one-sided mapping do not suffice for a Markov blocking kernel; the raw
  deterministic reflection-adapted average does preserve reflection positivity,
  and that distinction is a theorem.
\end{enumerate}

\section{Scope, standing hypotheses, and what would falsify this}\label{sec:scope}

\subsection{The claim boundary}

\begin{enumerate}\itemsep3pt
\item Everything in Sections~\ref{sec:orbits}--\ref{sec:su3} is exact rational
  arithmetic on a finite-dimensional projected operator at fixed lattice
  spacing. No limit is taken.
\item The engine's premises are unchanged by the $\Q(N)$ run and remain open
  routes: the Hermitian $PVP=0$ fourth-order form; the dropping of the $E_0$
  components (that the $E_0$ eigenspace on the reachable words is the plaquette
  span --- verified computationally on $224$ dropped components, not derived);
  the vacuum intermediates; and the completeness of the eleven-cumulant
  decomposition. The $\Q(N)$ run certifies that the spectral decomposition is
  exact and that the specialisation is faithful for $N\ge5$; it does not audit
  what the engine drops.
\item Lemma~\ref{lem:univ} is verified for $a,b\le2$, with $(2,2)$ sampled.
  The general-order statement is a conjecture.
\item The two named hypotheses gating the \emph{Hamiltonian} infinite-volume
  route in this program are the inhomogeneous Wilson free-energy bound with
  useful $\log K_\alpha$, and the source-radius reduction. Neither is proved
  here or anywhere in this program: the register carries them as its one
  load-bearing open gap, and the one attempt on them recorded to date fails on
  an $\alpha$-free constant and on a circular use of the gap it is meant to
  supply. The Euclidean counterpart of the needed uniform-in-volume
  strong-coupling control is classical constructive work; the open question is
  its Hamiltonian transcription for this projected sector.
\item A separate route \emph{does} now reach infinite volume at fixed spatial
  spacing, and this paper neither uses nor asserts it: an actual
  Perron-normalised Wilson transfer with a strong limit, a complete
  infinite-volume physical odd band, and a literal-source frame onto that band,
  on an explicit small-coupling interval. It is an analytic result with five
  stated hypotheses, at fixed spacing and fine temporal mesh, for $SU(N)$ with
  $N\ge3$; its interval and an asymptotically free trajectory are eventually
  disjoint domains. Nothing in Sections~\ref{sec:orbits}--\ref{sec:su3} depends
  on it, and it does not supply a continuum limit.
\item The continuum limit is open and two obstructions to the obvious route
  are theorems (Section~\ref{sec:notclaim}). What is missing is named rather
  than estimated: a Hamiltonian--Wilson matching package --- physical energy
  prefactor, anisotropy and speed-of-light renormalisation, the
  $g_H$--$g_W$ scheme map, and a renormalised spectral measure that can be
  followed to large $u$.
\item The charge-even identification remains a hypothesis. The unsigned
  incidence comparison operator has the cubic $\mu(\mu-p)^2=4\hat a_1\hat
  a_2\hat a_3$ with $p+q=12$ and range exactly $[-4,12]$; its $\Gamma$-point
  splitting is $12t_{r,+}$ \emph{conditional} on identifying it with the
  complete linked charge-even symbol, which is not derived.
\end{enumerate}

\subsection{Numerical archives that do not bear on these statements}

Because the identifiers used above ($W6$, $WR26$, $G17$, $G18$) name this
program's own open problems, one class of negative result belongs here rather
than in a footnote. A September~2026 GPU archive was supplied as resolving
several of them. Its arithmetic reproduces without a GPU and its deterministic
quantities have been replayed, so the numbers are not in question; what the
audit establishes is that the operators constructed are not the operators the
criteria are about. In summary, and each point is an identity rather than an
opinion:

\begin{itemize}\itemsep2pt
\item The ``fast-mode gauge resolvent'' is a commuting massive scalar model,
  $F_g=(1-g/4)K+(m^2+2g)I$ with $K=-\Delta$, containing no link variables, no
  gauge projection and no vacuum subtraction. It does admit a stronger result
  than the sampled scan --- a uniform bound $2g/[m^2(m^2+2g)]\le(32/81)g$ for
  all $g\ge0$, and a Combes--Thomas estimate
  $|R_g(x,y)|\le(2/m^2)e^{-\eta\,\mathrm{dist}}$ --- but not the operator
  identification the criterion requires.
\item The ``coupled domino'' is an exact tensor sum
  $h\otimes I+I\otimes h$ with no cross derivative; the omitted term is
  explicitly $\sum_a(L_a\chi)(U_1)(R_a\chi)(U_2)
  =\sin\theta_1\sin\theta_2\,n_1\!\cdot\!n_2$, orthogonal to all products of
  class functions, with squared Haar norm $3/16$.
\item The ``transfer matrix versus Langevin'' comparison diagonalises two
  radial discretisations and constructs no transfer matrix; the reported ratio
  is $0.93096\ldots$, and $\tau=t^2$ is not a semigroup clock, since
  $e^{-(s+t)^2L}\ne e^{-s^2L}e^{-t^2L}$ on a nonzero eigenmode.
\item The multiscale telescope is valid \emph{conditional on inserted inputs};
  its reported $\Delta_{20}$ is a finite iterate, and explicit two-sided bounds
  place the actual limit elsewhere.
\item The smearing statistic is a spatial low-momentum power fraction whose
  filter is independent of $a$; on the $L=8$ torus it satisfies
  $Z\le e^{-25\pi^2/1152}<0.85$ in exact arithmetic at every number of steps,
  so its asserted $\ge85\%$ uniform result is impossible for that statistic.
\item The Kotecký--Preiss part constructs no polymer weights; its own computed
  ratio at the largest quoted coupling is $52.475\ldots$, contradicting the
  stored verdict, and its ``free energy'' is a formula independent of $L$ for
  $L>8$.
\item The reflection-positivity matrix is $M=VV^{\mathsf T}$ by construction,
  hence positive semidefinite for \emph{every} input, and its zero eigenvalues
  are exact; the clustering part evaluates an inserted two-exponential formula
  and returns the inserted mass; the ``non-triviality'' statistic is the excess
  kurtosis of $\Tr U$ for a single group element, which is $-1$ exactly at
  $\beta=0$, i.e.\ nonzero for Haar measure with no interaction at all.
\item The eight accompanying Lean declarations compile and are registered under
  their literal statements. None concerns approximate tensorization, reflection
  positivity, or a four-point function; the one named for a bilinear pairing
  assumes $|R(j)|\le C\rho^j$ and concludes $\sum_{j<n}|R(j)|\le C/(1-\rho)$,
  an absolute scalar sum whose gauge parameter is unused.
\end{itemize}

None of this reflects on the numerical work as computation. It is recorded
because a criterion is discharged by identifying its operator, and a compiled
conditional statement is not its unconditional antecedent.

\subsection{The fourth order is not finished, and this is the part that is not}

Two things should be said explicitly, because a paper that closes a
long-standing dispute is exactly where a reader stops asking.

First, the \emph{$\Gamma$-scalar} side of the fourth order is not settled the
way the shape side is. The quantity $\Delta q_3=-16863189551/76406976000$ is a
\emph{definition} in this program, not a check: the direct $SU(3)$ fixed-rank
balanced contraction that would turn it into a measurement is absent, and a
Hadamard finite part misses it. Anything stated about $q_{\mathrm{band}}^{(4)}$
inherits that, as does anything containing the on-site amplitude $\tilde\sigma$
--- the one orbit amplitude that no independent lens has checked. The results
of this paper are about $A$, $B$, $D$ and $C_{\mathrm{shp}}$, none of which
contains $\tilde\sigma$.

Second, the register's own state as of this writing: no entry in the
contradiction register stands as open. Every recorded contradiction is
resolved, superseded or falsified. That is a statement about bookkeeping, not
about physics, and it is worth exactly as much as the checks behind each
verdict.

\subsection{Open problems, each with the experiment that settles it}

\begin{enumerate}\itemsep3pt
\item \textbf{Cumulant completeness at general $N$}
  (Hypothesis~\ref{hyp:complete}). Enumerate the fourth-order support at $N=5$
  from the engine and compare with the $189$ Hodge-form keys. Untried, and
  cheap.
\item \textbf{The $N=4$ pathway} (Hypothesis~\ref{hyp:resolvent}). Audit the
  determinant trick by an independent $\varepsilon$-network reduction, or
  restrict every headline to $N\ge5$.
\item \textbf{A second implementation at $N\ge4$} (Remark~\ref{rem:oneengine}).
  Recompute the corner dressing at $N=5$ by an independent route.
\item \textbf{The direct balanced contraction at $N=3$}
  (Remark~\ref{rem:adjudicated}). The declared falsifier of
  Theorem~\ref{thm:c2}.
\item \textbf{The universality lemma at full multiplicity}
  (Lemma~\ref{lem:univ}). Exhaust the $(2,2)$ swap family, or prove the lemma
  inductively in $a,b$; sixth order needs $a,b\le3$.
\item \textbf{The pair cluster} (Lemma~\ref{lem:pair}). Prove the graph
  isomorphism and the parity claim, or compute the cluster at $N=3$ and $N=4$.
\item \textbf{A cold regeneration of the corpus bundle}
  (Remark~\ref{rem:twopipelines}), which would turn output-certification into
  certification.
\item \textbf{The direct sixth-order assembly}
  (Proposition~\ref{prop:sixth}), which decides whether the $1:3:-4$ shape
  survives.
\item \textbf{Axiom guards on the headline Lean theorems}, and a recorded build
  with its toolchain version and exit code.
\item \textbf{The one external source with a recorded \emph{contradicts} edge}
  has not been obtained. It either becomes a measured disagreement or a
  convention difference that closes cleanly, and a referee will look for it
  first.
\end{enumerate}

\subsection{What would falsify the results of this paper}

\begin{itemize}\itemsep2pt
\item Theorem~\ref{thm:hodge}: a record of the fourth-order kernel not in the
  span of $\{I,\Ssq,\Ssq^2,\Lup,R\}$, or a residual in the Laurent solve.
\item Theorem~\ref{thm:support}: a word in the actual support whose carrier
  symbol leaves $\operatorname{span}\{q,q^2,e_2\}$.
\item Lemma~\ref{lem:univ}: one word of the family with
  $\int w\ne\int\varphi(w)$ --- most cheaply, inside the $99\%$ of the $(2,2)$
  swap family that has not been tested.
\item Theorem~\ref{thm:beta}: a rank at which a derived closed form disagrees
  with the per-rank engine, or a failure of the Lean identity. The identity is
  independently rerunnable in twenty lines of standard-library rational
  arithmetic from the coefficient strings printed in
  Appendices~\ref{app:closed} and~\ref{app:p17}; doing so reproduces
  $15644916262153/34416487661400$ at $N=3$ and
  $126537112003083861011/12716894720031723060840$ at $N=5$.
\item Theorem~\ref{thm:c2}: a twenty-fourth insertion ordering in the recovered
  ledger, or a fourth value for the corner cumulant from an implementation that
  shares no framework with this one.
\end{itemize}

\section{Evidence, tiers, and reproducibility}\label{sec:evidence}

This paper is written against a repository that decides, mechanically, which of
its statements are established. Four tiers are used, and an invisible marker in
this document's source names, after each statement it carries, the registered
check that establishes it.

\begin{center}
\begin{tabular}{@{}ll@{}}
\toprule
Tier & Meaning\\
\midrule
T0 & Lean~4 compiles it, no \texttt{sorry}, standard axioms only\\
T1 & re-derived symbolically from stated definitions, in exact rationals\\
T2 & float agreement within a tolerance printed in the check's detail line\\
T3 & a document says so and nothing checks it\\
\bottomrule
\end{tabular}
\end{center}

At the commit accompanying this paper the generated frontier reports
$562/562$ passing T1/T2 checks and $290$ Lean declarations with no
\texttt{sorry}. Three things about that sentence should be read carefully
rather than taken as a slogan.

\begin{enumerate}\itemsep3pt
\item \textbf{What T1 means here varies.} Some T1 checks recompute from
  definitions at verify time; many others read a pinned run certificate and
  compare it exactly. Both are exact, but the evidence in the second case is
  the pinned record plus its hash, not a fresh computation. The check's own
  detail line says which.
\item \textbf{The axiom claim is narrower than the tier label.} The build
  target is \texttt{lake build --wfail}, which detects \texttt{sorry}, not
  axioms. In-file \verb|#print axioms| guards cover $16$ of the $180$
  declarations, and the one recorded focused audit covered the $46$
  non-\texttt{Basic} declarations. The headline theorems of
  Section~\ref{sec:qn} live in \texttt{Basic.lean} and carry no such guard.
  Adding them is cheap and is listed among the open items below.
\item \textbf{No external group has reproduced anything in this program},
  and the internal replications share lineage: both cube builds read the same
  hash-pinned archive, and the from-scratch engine shares its \emph{framework}
  --- the Hodge form, the $8A+16C$ map, the normalisations --- with the
  pipeline it audits, because those were fixed against the corpus's own $N=3$
  records. Independence here is of values and code, not of framework, and a
  common-mode fault in the shared primitives would propagate identically into
  both.
\end{enumerate}

The suites
carrying this paper are: \textsf{fourth-order kernel orbits} ($28$),
\textsf{the third implementation and the historical ledger} ($11$),
\textsf{the third engine at every rank} ($12$), \textsf{the third engine over
$\Q(N)$} ($10$), \textsf{private-link paths; universality history by history}
($6$), \textsf{the Feshbach channel of the plaquette Hodge algebra} ($10$),
\textsf{the Feshbach resolvent comparison} ($5$), \textsf{off-axis channel
ledger} ($23$), \textsf{fourth order, sealed core} ($10$), \textsf{fourth
order, anchoring and the residual dispute} ($13$), \textsf{old-to-new
crosswalk} ($5$), \textsf{second order, all ranks} ($13$), \textsf{homology and
finite volume} ($15$), \textsf{the electric shell, and what isolates it} ($8$),
and \textsf{the Yang--Mills GPU archive: exact operator identification audit}
($6$).

Three principles govern how the numbers in this paper are to be read.

\begin{enumerate}\itemsep3pt
\item \textbf{Status and evidence are independent.} A claim can be proved in
  status and record-backed in evidence: the argument exists, the artifact does
  not. ``Certified'' is never a synonym for ``proved''.
\item \textbf{A failing check has found something.} In order: a bug in the
  check, a transcription slip in the registry, or a real discrepancy. For the
  third, the discrepancy is asserted by an explicit \textsc{finding} check ---
  which is why Findings~\ref{find:ur} and~\ref{find:rur} appear above as
  results rather than as caveats. A tolerance is never widened to make a
  finding disappear.
\item \textbf{Repetition is not corroboration.} With a corpus of this size and
  heavy copying, a value in forty files may have one origin; the join keys are
  exact rationals, not concepts.
\end{enumerate}

Each statement is re-established individually by the command printed with it in
the repository's certified list, and the whole exact layer runs in a few
seconds.

\section{Conclusion}\label{sec:conclusion}

The fourth order of the charge-odd flat band is now closed as an algebraic
object at every rank. Its kernel is the incidence algebra of the cubical
complex plus exactly one operator; the two top-degree shape coefficients vanish
because that algebra's carrier symbols span only $\{q,q^2,e_2\}$, with the one
dynamical input reduced to a finite, checked identity of Haar integrals; every
cluster cumulant is a derived rational function of $N$; their assembly is the
corpus's all-rank coefficient, as a polynomial identity proof-checked in Lean;
and the last open contradiction of the program --- the off-axis coefficient ---
is decided against both previously recorded kernels, with the error localised
to sixteen insertion orderings of one cluster.

What remains is not a coefficient. It is the two named hypotheses that gate
every infinite-volume statement, and the continuum matching package, whose
absence is now stated as a list of objects rather than as a difficulty. The
sixth order is the first place the machinery above can be tested rather than
extended: $\sigma(RR)=qe_2+3e_3$ says the $L^{-4}$ tier enters as soon as two
$R$'s can meet, at ratio $1:3$, and $\sigma(RUR)=4e_2^2$ says the shape ansatz
itself will not survive there.

\appendix

\section{The derived closed forms}\label{app:closed}

All in the C-odd sector unless stated. $z=N^2$ where convenient.

\paragraph{Two-hop weight.}
$u(N)=\mathcal N_u(N)/\mathcal D_u(N)$ with
\begin{align*}
  \mathcal N_u&=N^3(N^2-4)^2\bigl(16896N^{14}-131616N^{12}+451352N^{10}
   -882908N^{8}\\
   &\qquad\qquad\qquad +1058410N^{6}-771029N^{4}+313093N^{2}-54216\bigr),\\
  \mathcal D_u&=2(N^2-1)^3(4N^2-9)^3(9N^2-16)(2N^2-3)(2N^2-1)^3(3N^2-2)
   (4N^2-N-4)(4N^2+N-4).
\end{align*}
At $N=3$: $u=360421351/40327601932800$, and $u_2=2u$ exactly.

\paragraph{Single-contact dressing (perpendicular pair).}
\[
  d_{\mathrm{single}}(N)=\frac{2N^3(N^2-4)(10N^2-13)}
  {(N^2-1)^3(4N^2-9)^2(2N^2-1)^2}.
\]
In the C-even sector the same cluster is $-\ell_N^2/(2\CF)$, which is proved
in Lean.

\paragraph{Corner dressing.}
$d_{\mathrm{corner}}(N)=\mathcal N_c(N)/\mathcal D_c(N)$ with
\begin{align*}
  \mathcal N_c&=16N^3(N^2-4)\bigl(2379648N^{18}-28088736N^{16}+143075272N^{14}
    -411323454N^{12}\\
   &\quad +732994774N^{10}-837251963N^{8}+611821212N^{6}
    -275614672N^{4}+69464470N^{2}-7465875\bigr),\\
  \mathcal D_c&=(N^2-1)^3(4N^2-9)^3(4N^2-1)(9N^2-25)(2N^2-1)^3(3N^2-1)
    (4N^2-5)(4N^2-2N-5)(4N^2+2N-5).
\end{align*}
At $N=3$: $-2580244782961/398756546697600$ (C-odd) and
$-56022878647/4153714028100$ (C-even).

\paragraph{Cube completions.}
$K_{\mathrm{opposite}}=-160/[N(N^2-1)^3]$,
$K_{\mathrm{adjacent}}=-106/[N(N^2-1)^3]$, ratio $53/80$; and the axial cube
coefficient $\alpha_N=640/[N(N^2-1)^3]$, with $\alpha_3=5/12$.

\paragraph{Second order, for reference.}
$t_N=2N(N^2-4)/[(N^2-1)(2N^2-1)(4N^2-9)]$, $t_3=5/612$;
$\CF=(N^2-1)/(2N)$; the four shared-link channel weights are $d_\rho/N^2$; and
the planar closed form
$1-4N^3t_N=(2N^4+31N^2-9)/[(N^2-1)(2N^2-1)(4N^2-9)]$.

\section{$P_{17}$ and $R_{20}$}\label{app:p17}

With $z=N^2$,
\begin{align*}
 P_{17}(z)&=2096187310080\,z^{17}-45206560309248\,z^{16}+448972002607104\,z^{15}\\
 &\quad-2723575470882816\,z^{14}+11288692151812096\,z^{13}-33888218411529728\,z^{12}\\
 &\quad+76218901019673664\,z^{11}-131068691814847264\,z^{10}+174326341061538992\,z^{9}\\
 &\quad-180230597250871976\,z^{8}+144751635142984472\,z^{7}-89742150515602808\,z^{6}\\
 &\quad+42388925672412712\,z^{5}-14916377727371552\,z^{4}+3768794520714128\,z^{3}\\
 &\quad-641987460459360\,z^{2}+65414604672000\,z-2967321600000,
\end{align*}
\begin{align*}
 R_{20}(z)&=(z-1)^3(2z-3)(2z-1)^3(3z-2)(3z-1)(4z-9)^3(4z-5)(4z-1)\\
 &\qquad\times(9z-25)(9z-16)(16z^2-44z+25)(16z^2-33z+16).
\end{align*}
Each cumulant denominator divides $NR_{20}(N^2)$; the four cofactors are given
explicitly in the Lean source, and $\beta_3=15644916262153/34416487661400$.

\section{The $SU(3)$ orbit ledger}\label{app:ledger}

\begin{center}
\begin{tabular}{@{}llr@{}}
\toprule
Orbit & Records & Amplitude at $N=3$\\
\midrule
skeleton ($u$) & $132$ & $360421351/40327601932800$\\
same-plane $(0,1,1)$ in-plane ($u_2$) & $12$ & $360421351/20163800966400=2u$\\
cross-plane rotation ($\rho$) & $24$ & $238714892212171339/29002361154409843200$\\
coplanar ($\tilde\pi$) & $12$ & $-20535103905179/1264270320593280$\\
normal ($\tilde\nu$) & --- & $-5/48$\\
on-site ($\tilde\sigma$) & --- & $-780864191400383617/302107928691769200$\\
\midrule
\multicolumn{2}{@{}l}{$C_{\mathrm{shp}}$, historical} & $-211835444920651/4405310420659200$\\
\multicolumn{2}{@{}l}{$C_{\mathrm{shp}}$, assembled (this paper)} & $-13035490122347/550663802582400$\\
\bottomrule
\end{tabular}
\end{center}

\section{Machine-check coverage}\label{app:coverage}

The repository generates, from the registered checks themselves, the list of
every marker in this paper together with the check it names
and that check's live verdict. A marker naming a check that does not exist, or
that does not pass, fails a test. The coverage list is not reproduced here; it
is generated at build time and is the authoritative statement of which
sentences above are carried by a machine check and which are prose.

\end{document}
